\documentclass[11pt,a4paper]{article}
\usepackage[utf8]{inputenc}
\usepackage[T1]{fontenc}
\usepackage[italian]{babel}
\usepackage{amsmath,amssymb,mathtools}
\usepackage[top=2.2cm,bottom=2.2cm,left=2.3cm,right=2.3cm]{geometry}
\usepackage{enumitem}
\usepackage{array}
\usepackage{booktabs}
\usepackage{titlesec}
\usepackage[most]{tcolorbox}
\usepackage{tikz}
\usetikzlibrary{arrows.meta,positioning,calc}
\usepackage{hyperref}

\definecolor{navy}{HTML}{1F3A63}
\definecolor{teal}{HTML}{0F6E6E}
\definecolor{amber}{HTML}{B26B00}
\definecolor{crimson}{HTML}{9B2226}
\definecolor{lightbg}{HTML}{F6F7FA}

\hypersetup{colorlinks=true,linkcolor=navy,urlcolor=navy}

% --- spaziatura generosa ---
\setlength{\parindent}{0pt}
\setlength{\parskip}{0.55em}
\linespread{1.08}
\setlist[itemize]{itemsep=0.35em,topsep=0.35em,leftmargin=1.4em}
\setlist[enumerate]{itemsep=0.4em,topsep=0.35em,leftmargin=1.6em}

% --- titoli ---
\titleformat{\section}{\Large\bfseries\color{navy}}{\thesection.}{0.6em}{}[\vspace{-0.5em}\color{navy}\rule{\textwidth}{1pt}]
\titlespacing*{\section}{0pt}{1.6em}{1em}
\titleformat{\subsection}{\large\bfseries\color{teal}}{\thesubsection}{0.6em}{}
\titlespacing*{\subsection}{0pt}{1.4em}{0.6em}
\titleformat{\subsubsection}{\bfseries\color{navy}}{\thesubsubsection}{0.6em}{}
\titlespacing*{\subsubsection}{0pt}{1.1em}{0.4em}

% --- box ---
\newtcolorbox{ric}[1]{enhanced,breakable,colback=teal!4,colframe=teal,boxrule=1pt,arc=4pt,
  left=10pt,right=10pt,top=8pt,bottom=8pt,fonttitle=\bfseries,coltitle=white,
  title={Ricetta \ --- \ #1}}
\newtcolorbox{att}[1]{enhanced,breakable,colback=crimson!4,colframe=crimson,boxrule=1pt,arc=4pt,
  left=10pt,right=10pt,top=8pt,bottom=8pt,fonttitle=\bfseries,coltitle=white,
  title={Attenzione \ --- \ #1}}
\newtcolorbox{teo}[1]{enhanced,breakable,colback=navy!4,colframe=navy,boxrule=1pt,arc=4pt,
  left=10pt,right=10pt,top=8pt,bottom=8pt,fonttitle=\bfseries,coltitle=white,
  title={Teoria \ --- \ #1}}
\newtcolorbox{es}[1]{enhanced,breakable,colback=amber!5,colframe=amber,boxrule=1pt,arc=4pt,
  left=10pt,right=10pt,top=8pt,bottom=8pt,fonttitle=\bfseries,coltitle=white,
  title={Esempio svolto \ --- \ #1}}
\newtcolorbox{idea}{enhanced,breakable,colback=amber!6,colframe=amber,boxrule=0pt,
  leftrule=4pt,arc=0pt,left=10pt,right=10pt,top=7pt,bottom=7pt}

% blocco "mattone": nome + formula, ben distanziati
\newcommand{\mattone}[2]{%
\par\vspace{0.5em}\textbf{#1}\par\vspace{-0.4em}
\begin{center}$\displaystyle #2$\end{center}\vspace{-0.2em}}


% --- grafi per l'esempio di cancellazione dei cicli ---
\tikzset{nd/.style={circle,draw=navy,thick,minimum size=7mm,inner sep=1pt,font=\small}}
\newcommand{\mcfnodes}{%
  \node[nd] (A) at (0,2.1) {$A$};
  \node[nd] (B) at (0,-2.1) {$B$};
  \node[nd] (C) at (3.0,0) {$C$};
  \node[nd] (D) at (6.0,2.1) {$D$};
  \node[nd] (E) at (6.0,-2.1) {$E$};}
% arco residuo: #1 = tipo di freccia, #2 = da, #3 = a, #4 = costo diretto
\newcommand{\ra}[4]{\draw[#1,thick,black!65] (#2) -- node[font=\scriptsize,black!70,fill=amber!5,inner sep=1pt]{$#4$} (#3);}
% arco del ciclo scelto
\newcommand{\cy}[3]{\draw[->,line width=1.2pt,crimson] (#1) to[bend left=16] node[font=\scriptsize,crimson,fill=amber!5,inner sep=1pt]{$#3$} (#2);}
% arco del grafo originale: #3 = (capacita,costo), #4 = flusso
\newcommand{\fa}[4]{\draw[->,thick,black!45] (#1) -- node[font=\scriptsize,black!55,sloped,above,inner sep=1.5pt]{$#3$} node[font=\small,teal,sloped,below,inner sep=1.5pt]{$\mathbf{#4}$} (#2);}

\begin{document}

%---------------------------------------------------------------- copertina
\thispagestyle{empty}
\vspace*{3cm}
\begin{center}
{\color{navy}\rule{\textwidth}{2pt}}\\[1.2em]
{\Huge\bfseries\color{navy} Formulario di\\[0.3em] Ottimizzazione Combinatoria}\\[1.2em]
{\large Spiegazioni passo passo ed esempi svolti dai compiti d'esame}\\[1em]
{\color{navy}\rule{\textwidth}{2pt}}
\end{center}
\vfill
\begin{idea}
\textbf{Come si usa.} Ogni sezione copre un esercizio dello scritto.

\textbf{Ricette} (verde): la procedura da seguire meccanicamente durante il compito.\\
\textbf{Esempi svolti} (arancione): lo stesso procedimento applicato a un compito vero, con tutti i conti.\\
\textbf{Teoria} (blu): enunciati e dimostrazioni da riportare tali e quali.\\
\textbf{Attenzione} (rosso): gli errori che tolgono punti.

Gli esempi svolti sono presi dai compiti del 25/05/2026, 25/06/2026, 27/07/2026 e 08/07/2025
e i conti sono stati rifatti a mano: puoi fidarti dei numeri e usarli come autoverifica.
\end{idea}
\newpage

\thispagestyle{empty}
\tableofcontents
\newpage

%================================================================
\section{Struttura dello scritto}

Tre esercizi, sempre degli stessi tipi.

\begin{center}
\begin{tabular}{@{}p{0.9cm} p{5.0cm} p{7.9cm}@{}}
\toprule
& \textbf{Argomento} & \textbf{Cosa viene chiesto}\\
\midrule
\textbf{Es 1} & Modellazione PLI \newline (8--9 punti) &
Testo a parole $\to$ variabili, obiettivo, vincoli. A volte una variante del testo (punti extra).\\
\addlinespace
\textbf{Es 2} & PL e simplesso \newline (11--13 punti) &
(a) regione ammissibile grafica, (b) simplesso primale da una base data, (c) duale oppure una definizione, (d) una dimostrazione.\\
\addlinespace
\textbf{Es 3} & Reti di flusso \newline (8--11 punti) &
Edmonds--Karp, oppure costo minimo (cancellazione cicli), oppure matching; più una o due domande di teoria.\\
\bottomrule
\end{tabular}
\end{center}

\begin{idea}
\textbf{Regola d'oro.} Scrivi tutti i passaggi. I punti si prendono anche col procedimento
corretto e il conto sbagliato; non si prendono mai col risultato giusto e nessuna spiegazione.
Il testo dice esplicitamente: \emph{«occorre riportare in modo chiaro tutti i passi che portano
alla determinazione del risultato»}.
\end{idea}

\begin{att}{niente calcolatrice}
Non sono ammesse calcolatrici. Tutti i conti del simplesso sono su matrici $2\times 2$
con numeri piccoli: si fanno a mano con la formula dell'inversa (sezione~\ref{sec:inversa}).
Se ti escono frazioni mostruose, quasi certamente hai sbagliato un segno nella normalizzazione.
\end{att}

\newpage
%================================================================
\section{Esercizio 1 --- Modellazione PLI}

\subsection{La procedura}

\begin{ric}{modellazione in cinque passi}
\begin{enumerate}
\item \textbf{Dati.} Elenca insiemi e parametri con i loro indici. Scriverli \emph{prima}
      evita di inventare simboli a metà esercizio.
      Tipicamente: $i = 1,\dots,m$ \emph{contenitori} (macchine, camion, camerieri, server),
      $j = 1,\dots,n$ \emph{oggetti} (lavori, pallet, tavoli, query).

\item \textbf{Variabili.} Decidi \emph{cosa} si sceglie. Binarie per sì/no, intere per quantità.
      Scrivi sempre il significato a parole, altrimenti si perdono punti anche col modello giusto:
      \[ x_{ij} = \begin{cases} 1 & \text{se l'oggetto $j$ è assegnato a $i$}\\ 0 & \text{altrimenti}\end{cases}
         \qquad
         y_{i} = \begin{cases} 1 & \text{se $i$ è attivato}\\ 0 & \text{altrimenti}\end{cases}\]
      \emph{Regola pratica:} se nel testo compare un \textbf{costo fisso} (di utilizzo, di
      apertura, di assunzione) serve \emph{sempre} anche la variabile $y_i$.

\item \textbf{Obiettivo.} Somma di costi o profitti moltiplicati per le variabili.
      Controlla se è minimo o massimo. Se ci sono costi fissi \emph{e} variabili,
      l'obiettivo ha due sommatorie: $\sum_i (\text{fisso})_i\, y_i + \sum_i\sum_j (\text{variabile})_{ij}\, x_{ij}$.

\item \textbf{Vincoli.} Uno per ogni frase del testo. Rileggi il testo frase per frase
      e per ognuna scrivi la disuguaglianza corrispondente. Quantifica sempre con $\forall i$ oppure $\forall j$.

\item \textbf{Domini.} $x_{ij} \in \{0,1\}$, \quad $y_i \in \{0,1\}$, \quad $z_i \ge 0$ intero.
\end{enumerate}
\end{ric}

\begin{idea}
\textbf{Come si capisce quale indice va sotto la sommatoria.} Guarda \emph{su cosa} il vincolo
parla:
\begin{itemize}
\item «ogni \emph{oggetto} $j$ deve\dots» $\Rightarrow$ somma su $i$ (i contenitori), vincolo $\forall j$;
\item «ogni \emph{contenitore} $i$ può\dots» $\Rightarrow$ somma su $j$ (gli oggetti), vincolo $\forall i$.
\end{itemize}
L'indice che compare nel \emph{«per ogni»} è quello che \emph{non} viene sommato.
\end{idea}

\subsection{I mattoni da memorizzare}

Quasi ogni testo d'esame è una combinazione di questi blocchi.

\mattone{Assegnamento --- ogni $j$ va a esattamente una risorsa}
{\sum_{i} x_{ij} = 1 \qquad \forall j}

\mattone{Semi-assegnamento --- ogni $j$ va al più a una risorsa (\emph{«un sottoinsieme di\dots»})}
{\sum_{i} x_{ij} \le 1 \qquad \forall j}

\mattone{Copertura multipla --- ogni $j$ va su esattamente $p_j$ risorse}
{\sum_{i} x_{ij} = p_j \qquad \forall j}

\mattone{Capacità --- peso o tempo massimo della risorsa $i$}
{\sum_{j} w_j\, x_{ij} \le b_i \qquad \forall i}

\mattone{Cardinalità --- al più $q_i$ oggetti sulla risorsa $i$}
{\sum_{j} x_{ij} \le q_i \qquad \forall i}

\mattone{Attivazione e costo fisso --- lega $x$ alla risorsa aperta $y$ (due forme equivalenti)}
{\sum_{j} w_j x_{ij} \le b_i\, y_i \quad \forall i
 \qquad\text{oppure}\qquad
 \Big[\ \sum_{j} w_j x_{ij} \le b_i \ \ \forall i \ \ \text{ e } \ \ x_{ij} \le y_i \ \ \forall i,j\ \Big]}

\mattone{Incompatibilità --- $h$ e $k$ non sulla stessa risorsa}
{x_{ih} + x_{ik} \le 1 \qquad \forall i, \ (h,k) \in L}

\mattone{Prerequisito --- $b$ richiede $a$}
{x_b \le x_a}

\mattone{Quorum --- almeno metà delle risorse attive}
{\sum_{i} y_i \ \ge\ \left\lceil n/2 \right\rceil}

\mattone{Quantità minima positiva --- o $x = 0$, oppure $x \ge \ell$}
{\ell\, y \ \le\ x \ \le\ M\, y, \qquad y \in \{0,1\}}

\mattone{Big-M --- vincolo attivo solo quando $y = 1$}
{a^T x \ \le\ b + M\,(1 - y)}

\mattone{Valore assoluto --- linearizzare $\min |z|$}
{\min t \quad\text{con}\quad t \ge z, \ \ t \ge -z}

\mattone{Min-max, tipo makespan --- minimizzare il carico più alto}
{\min T \quad\text{con}\quad T \ \ge\ \sum_{j} p_j\, x_{ij} \quad \forall i}

\newpage
\subsection{Un testo tradotto frase per frase}

\begin{es}{camion e pallet (compito 25/05/2026)}
\textbf{Testo.} \emph{«Un'azienda deve caricare gli $m$ camion disponibili per servire un
sottoinsieme di clienti, ciascuno dei quali richiede un pallet. Ogni camion $i$ ha capacità
$b_i$ quintali e costo di utilizzo $c_i$ euro. Si scelgono un sottoinsieme degli $n$ pallet,
ciascuno con profitto $p_j$ e peso $w_j$. Massimizzare il profitto netto. Qualche camion può
restare inutilizzato a costo nullo.»}

\medskip
\textbf{Traduzione, frase per frase:}

\begin{center}
\begin{tabular}{@{}p{7.0cm} p{7.4cm}@{}}
\toprule
\textbf{Frase del testo} & \textbf{Pezzo di modello}\\
\midrule
«caricare gli $m$ camion» & serve $x_{ij}=1$ se il pallet $j$ va sul camion $i$\\
\addlinespace
«costo di utilizzo $c_i$», «qualche camion può restare inutilizzato» &
serve $y_i=1$ se il camion $i$ è usato: il costo si paga solo se lo si usa\\
\addlinespace
«un \emph{sottoinsieme} dei pallet» &
semi-assegnamento: $\sum_i x_{ij} \le 1$, non $=1$\\
\addlinespace
«capacità in peso $b_i$» &
$\sum_j w_j x_{ij} \le b_i y_i$: con $y_i=0$ il camion non può portare nulla\\
\addlinespace
«massimizzare il profitto netto = profitto pallet $-$ costo camion» &
$\max \sum_i\sum_j p_j x_{ij} - \sum_i c_i y_i$\\
\bottomrule
\end{tabular}
\end{center}

\medskip
\textbf{Modello finale.}
\[
\max \ \sum_{i=1}^{m}\sum_{j=1}^{n} p_j x_{ij} \ -\ \sum_{i=1}^{m} c_i y_i
\]
\[
\begin{aligned}
&\sum_{i=1}^{m} x_{ij} \le 1 && \forall j = 1,\dots,n &&\text{(ogni pallet al più su un camion)}\\
&\sum_{j=1}^{n} w_j x_{ij} \le b_i\, y_i && \forall i = 1,\dots,m &&\text{(capacità + attivazione)}\\
&x_{ij} \in \{0,1\} && \forall i,j\\
&y_i \in \{0,1\} && \forall i
\end{aligned}
\]

\textbf{Perché $b_i y_i$ e non $b_i$.} Se scrivessi $\sum_j w_j x_{ij} \le b_i$, allora
$y_i$ non comparirebbe in nessun vincolo: l'ottimo metterebbe $y_i = 0$ per tutti (risparmiando
i costi $c_i$) e caricherebbe comunque i pallet. Il modello sarebbe sbagliato.
La forma alternativa accettata è $\sum_j w_j x_{ij}\le b_i$ \emph{insieme a} $x_{ij}\le y_i$
$\forall i,j$: anche lì $y_i$ è forzato a 1 appena un pallet è caricato.
\end{es}

\newpage
\subsection{Archetipi già usciti all'esame}

\textbf{Camerieri e tavoli} \quad {\small(compito 25/06/2026)}\\
Costo fisso $l_i$ per cameriere assunto, costo variabile $v_i$ per tavolo servito, tempo $p_j$
per servire il tavolo $j$, budget di tempo $c_i$ per cameriere. Qui i tavoli vanno serviti
\emph{tutti}, quindi assegnamento con $=1$.
\[
\min \ \sum_i l_i y_i + \sum_i \sum_j v_i x_{ij}
\qquad
\begin{aligned}
&\sum_i x_{ij} = 1 &&\forall j\\
&\sum_j p_j x_{ij} \le c_i\, y_i &&\forall i\\
&x_{ij},\,y_i \in \{0,1\}
\end{aligned}
\]

\textbf{Lavori e macchine con lista di incompatibilità} \quad {\small(compito 27/07/2026)}\\
Ogni lavoro $j$ va eseguito da $p_j$ macchine, ogni macchina ne fa al più $q_i$,
e la lista $L$ contiene coppie che non possono stare sulla stessa macchina.
Qui \emph{non} c'è costo fisso, quindi \emph{non} servono le $y_i$.
\[
\min \ \sum_i \sum_j c_{ij} x_{ij}
\qquad
\begin{aligned}
&\sum_i x_{ij} = p_j &&\forall j\\
&\sum_j x_{ij} \le q_i &&\forall i\\
&x_{ih} + x_{ik} \le 1 &&\forall i,\ \forall (h,k) \in L\\
&x_{ij} \in \{0,1\}
\end{aligned}
\]

\textbf{Server e query} \quad {\small(compito 08/07/2025)}\\
Attenzione: qui l'oggetto non è individuale, sono $k$ query \emph{indistinguibili}.
Quindi la variabile è \textbf{intera}, non binaria: $x_i$ = numero di query assegnate al server $i$.
Il testo chiede di tenere accesi almeno metà degli $n$ server.
\[
\min \ \sum_i p_i x_i
\qquad
\begin{aligned}
&\sum_i x_i = k && \text{(tutte le query vanno smaltite)}\\
&x_i \le q_i\, y_i &&\forall i \quad\text{(capacità + accensione)}\\
&\sum_i y_i \ge \lceil n/2 \rceil && \text{(vincolo contrattuale)}\\
&x_i \ge 0 \ \text{intero},\quad y_i \in \{0,1\}
\end{aligned}
\]

\textbf{Variante dello stesso compito} (punti extra): togliere il vincolo di quorum e
aggiungere un consumo $a_i$ per il solo fatto che il server $i$ resti acceso. Cambia solo
l'obiettivo, i vincoli restano (meno il quorum):
\[
\min \ \sum_i p_i x_i + \sum_i a_i y_i
\]

\begin{idea}
\textbf{Riconoscere il tipo in dieci secondi.}
Costo fisso nel testo $\Rightarrow$ variabili $y$. \enspace
«tutti» / «ogni\dots deve» $\Rightarrow$ vincolo $=$. \enspace
«un sottoinsieme» / «se conveniente» $\Rightarrow$ vincolo $\le$. \enspace
Oggetti indistinguibili e contati $\Rightarrow$ variabili intere, non binarie.
\end{idea}

\begin{att}{errori che costano punti}
\begin{itemize}
\item Vincolo di capacità scritto senza $y_i$ quando c'è un costo fisso: il costo non viene mai
      pagato e il modello è sbagliato (è l'errore più penalizzato).
\item Quantificatori $\forall i$ o $\forall j$ dimenticati.
\item Variabili continue dove servono binarie, o binarie dove serve una quantità intera.
\item Obiettivo che mescola minimizzazione e massimizzazione.
\item Variabili non definite a parole.
\item Dimenticare la riga finale dei domini.
\end{itemize}
\end{att}

\newpage
%================================================================
\section{Esercizio 2 --- PL, dualità, simplesso}

\subsection{Forma standard del corso}

\begin{att}{tutto va portato in questa forma, sempre}
Il corso lavora \emph{esclusivamente} con
\[
\boxed{\ \max\ c\,x \qquad A x \le b\ }
\]
$A$ è $m\times n$ (all'esame $n=2$ variabili, $m=6$ vincoli), $c$ è un vettore \emph{riga},
$x$ e $b$ vettori colonna. \textbf{Non} ci sono vincoli di non negatività impliciti:
se il testo vuole $x_1\ge 0$, quello è un vincolo come gli altri e va scritto come $-x_1 \le 0$.
\end{att}

\begin{ric}{normalizzazione}
\begin{enumerate}
\item \textbf{Obiettivo.} Se il testo dice $\min c\,x$, riscrivi $\max\, (-c)\,x$.
      Ricordati alla fine che il valore ottimo del problema originale è $-$(valore trovato).
\item \textbf{Vincoli $\ge$.} Moltiplica per $-1$ e girali: $a\,x \ge b \Longrightarrow -a\,x \le -b$.
\item \textbf{Vincoli $=$.} Si sdoppiano in $a x \le b$ e $-a x \le -b$ (raro all'esame).
\item \textbf{Variabili a sinistra, costanti a destra.} Porta ogni $x$ a sinistra del segno
      e ogni numero a destra.
\item \textbf{Numera i vincoli} $1,\dots,m$ e scrivi esplicitamente $A$, $b$, $c$.
\end{enumerate}
\end{ric}

\subsubsection{Come si numerano i vincoli e come si individua la base iniziale}
\label{sec:base}

Il testo d'esame presenta i vincoli in una tabella su \textbf{tre colonne da due righe},
e poi dice: \emph{«partendo dalla base costituita dai due vincoli della colonna più a sinistra»}.

\begin{idea}
\textbf{Regola di numerazione:} si legge \textbf{colonna per colonna, dall'alto in basso}.
Quindi la prima colonna dà i vincoli $1$ e $2$, la seconda i vincoli $3$ e $4$, la terza i
vincoli $5$ e $6$. La base iniziale è di conseguenza sempre
\[ B = \{1,\,2\}. \]
\end{idea}

\textbf{Che cos'è concretamente una base, in questo corso.} Non è un insieme di variabili
(come nel simplesso a tableau di altri corsi): è un insieme di \textbf{indici di vincoli}.
Con $n=2$ variabili, una base è una coppia di vincoli. Dalla base si ricava tutto:

\begin{center}
\begin{tabular}{@{}p{3.2cm} p{11.3cm}@{}}
\toprule
\textbf{Simbolo} & \textbf{Che cos'è, in pratica}\\
\midrule
$B$ & i due indici di riga scelti, es.\ $B=\{1,2\}$\\
$A_B$ & la matrice $2\times 2$ fatta \emph{solo} da quelle due righe di $A$\\
$b_B$ & il vettore $2\times 1$ con i due termini noti corrispondenti\\
$N$ & tutti gli altri indici, es.\ $N=\{3,4,5,6\}$\\
$A_N,\ b_N$ & le righe restanti di $A$ e le componenti restanti di $b$\\
$x$ & $A_B^{-1} b_B$: il punto dove i due vincoli di base valgono con l'\emph{uguaglianza}\\
\bottomrule
\end{tabular}
\end{center}

Geometricamente: i due vincoli di base sono due rette; la soluzione di base è il loro
punto di intersezione. È una \emph{base} legittima solo se le due rette non sono parallele,
cioè se $\det A_B \ne 0$. È \emph{ammissibile} se quel punto soddisfa anche gli altri quattro
vincoli, e in tal caso è un vertice della regione ammissibile.

\begin{att}{la base data dal testo è già ammissibile}
All'esame la base iniziale $\{1,2\}$ è sempre ammissibile: non serve una fase~I.
Se calcolando $x = A_B^{-1}b_B$ ti accorgi che qualche vincolo fuori base è violato,
hai sbagliato la normalizzazione (tipicamente un segno). Torna indietro e ricontrolla.
\end{att}

\subsubsection{Inversa di una matrice $2\times 2$}
\label{sec:inversa}

Tutta l'algebra del simplesso all'esame si riduce a questa formula. Impararla bene fa
risparmiare metà del tempo.
\[
A_B = \begin{pmatrix} a & b\\ c & d\end{pmatrix}
\qquad
\det A_B = ad - bc
\qquad
A_B^{-1} = \frac{1}{ad-bc}\begin{pmatrix} d & -b\\ -c & a\end{pmatrix}
\]
A parole: \emph{scambia i due elementi della diagonale principale, cambia segno agli altri due,
dividi tutto per il determinante}.

\textbf{Controllo veloce:} $A_B A_B^{-1}$ deve dare $\begin{psmallmatrix}1&0\\0&1\end{psmallmatrix}$.
Costa cinque secondi e salva l'intera iterazione.

\subsection{Metodo grafico}

Sono tre punti facili: vale sempre la pena farli con cura, e servono anche a controllare
il risultato del simplesso.

\begin{ric}{risoluzione grafica}
\begin{enumerate}
\item Normalizza i vincoli e numerali (sezione~\ref{sec:base}).
\item Per ogni vincolo $a_1 x_1 + a_2 x_2 \le b$, disegna la \textbf{retta} $a_1x_1+a_2x_2 = b$.
      Il modo più veloce: trova i due punti in cui incrocia gli assi
      (poni $x_1=0$ e ricava $x_2$, poi $x_2=0$ e ricava $x_1$) e uniscili.
      Se il vincolo ha una sola variabile è una retta verticale o orizzontale.
\item Per capire da che lato sta il semipiano ammissibile, \textbf{sostituisci un punto di prova}:
      se l'origine non sta sulla retta, usa $(0,0)$. Se $a_1\cdot 0 + a_2\cdot 0 \le b$ è vera,
      il semipiano buono è quello che contiene l'origine; altrimenti è l'altro.
      Segna una freccetta sulla retta verso il lato buono.
\item Colora l'intersezione dei sei semipiani: è la regione ammissibile. Segna i vertici con un pallino.
\item Disegna il vettore gradiente $c$ uscente dall'origine (per $c=(1,1)$: la bisettrice
      verso l'alto a destra). Essendo il problema in forma $\max$, ci si muove \textbf{nel verso di $c$}.
\item Il vertice più spinto in quel verso è l'ottimo. Ricavane le coordinate risolvendo il sistema
      dei due vincoli attivi lì (quello è anche il modo di leggere la base ottima!), poi valuta $c\,x$.
\end{enumerate}
\end{ric}

\begin{idea}
\textbf{Il legame fra grafico e simplesso.} Ogni vertice del disegno corrisponde a una base;
i vincoli che passano per quel vertice sono gli indici in $B$. Se il simplesso ti dà una base
finale $B=\{3,5\}$, sul disegno devi vedere le rette $3$ e $5$ incrociarsi proprio nel vertice
ottimo. È l'autoverifica migliore che hai.
\end{idea}

\newpage
\subsection{Simplesso primale: la ricetta del corso}

\begin{ric}{una iterazione (problema $\max c\,x$, $Ax \le b$)}
Dati: base $B$ (due indici), $A_B$, $b_B$, e $N$ con $A_N$, $b_N$.

\begin{enumerate}
\item \textbf{Inversa e punto corrente.}
      \[ A_B^{-1} \quad\text{(formula }2\times2\text{)}, \qquad x = A_B^{-1}\, b_B \]
      Alla prima iterazione verifica che $x$ soddisfi anche i vincoli fuori base.

\item \textbf{Soluzione duale di base (costi ridotti).}
      \[ y_B = c\, A_B^{-1} \]
      Attenzione all'ordine: $c$ è un vettore \emph{riga} e moltiplica $A_B^{-1}$ \emph{da sinistra}.
      Il risultato è un vettore riga con una componente per ogni vincolo in base.
      Il vettore duale completo $y$ ha quelle componenti sugli indici di $B$ e zero altrove.

\item \textbf{Test di ottimalità.} Se $y_B \ge 0$ (tutte le componenti), \textbf{stop}:
      $x$ è ottimo, e $y$ è il certificato duale. Valore ottimo $c\,x = y\,b$.

\item \textbf{Vincolo uscente.} Scegli $h \in B$ con $y_h < 0$.
      Se ce n'è più di uno: regola di Bland, prendi l'indice più piccolo (evita i cicli).

\item \textbf{Direzione.}
      \[ \xi = -\,A_B^{-1}\, u_h \]
      dove $u_h$ è il versore che vale 1 nella \emph{posizione che $h$ occupa dentro $B$}
      (prima o seconda). In pratica: \textbf{prendi quella colonna di $A_B^{-1}$ e cambiale segno}.
      Questa direzione lascia attivo l'altro vincolo di base e allontana da $h$, aumentando l'obiettivo.

\item \textbf{Rapporti.} Calcola i due vettori $A_N \xi$ e $A_N x$. Per ogni $i \in N$ con
      $A_i \xi > 0$ (cioè i vincoli verso cui ci stiamo muovendo):
      \[ \lambda_i \;=\; \frac{b_i - A_i x}{A_i \xi} \;\ge\; 0 \]
      cioè \emph{slack diviso velocità di avvicinamento}.
      Se \textbf{nessun} $i$ ha $A_i\xi > 0$, il problema è \textbf{illimitato}: stop.

\item \textbf{Vincolo entrante.} $k = \arg\min_{i} \lambda_i$ --- ovvero l'\textbf{indice} con $\lambda$ minimo.
      Nuovo punto e nuova base:
      \[ x' = x + \lambda_k\, \xi, \qquad B' = \big(B \setminus \{h\}\big) \cup \{k\} \]

\item Torna al passo 1. (Conviene riordinare $B$ in ordine crescente e ricalcolare $N$.)
\end{enumerate}
\end{ric}

\begin{att}{i tre errori tipici}
\begin{itemize}
\item \textbf{$y_B = c A_B^{-1}$, non $A_B^{-1} c$.} Il prodotto è riga per matrice.
\item \textbf{Il segno di $\xi$.} C'è un meno davanti: $\xi = -A_B^{-1}u_h$.
\item \textbf{Il verso del rapporto.} Si prendono solo i vincoli con $A_i\xi > 0$
      (in forma $\le$ ci si avvicina al vincolo quando $A_i\xi$ è \emph{positivo}).
      È l'opposto di quello che si ricorda dai libri scritti in forma $\ge$.
\end{itemize}
\end{att}

\newpage
\subsection{Esempio svolto completo di simplesso}

\begin{es}{compito del 27/07/2026, esercizio 2}

\textbf{Testo.}
\[ \max\ x_1 + x_2 \]
\begin{center}
\begin{tabular}{@{}ccc@{}}
$x_2 + 2 \ge -x_1$ & $x_1 \le 2$ & $x_1 \le 4 - x_2$\\
$x_2 \ge -1$ & $x_2 \ge x_1 - 1$ & $2x_1 \ge 3x_2 - 6$
\end{tabular}
\end{center}
Base iniziale: i due vincoli della colonna più a sinistra.

\medskip
\hrule
\medskip

\textbf{Passo 0 --- normalizzazione.} Leggendo per colonne, dall'alto in basso:

\begin{center}
\begin{tabular}{@{}c l l@{}}
\toprule
$i$ & \textbf{come è scritto} & \textbf{normalizzato ($\le$)}\\
\midrule
1 & $x_2 + 2 \ge -x_1$ & $-x_1 - x_2 \le 2$\\
2 & $x_2 \ge -1$ & $-x_2 \le 1$\\
3 & $x_1 \le 2$ & $x_1 \le 2$\\
4 & $x_2 \ge x_1 - 1$ & $x_1 - x_2 \le 1$\\
5 & $x_1 \le 4 - x_2$ & $x_1 + x_2 \le 4$\\
6 & $2x_1 \ge 3x_2 - 6$ & $-2x_1 + 3x_2 \le 6$\\
\bottomrule
\end{tabular}
\end{center}

\[
A = \begin{pmatrix}
-1 & -1\\ 0 & -1\\ 1 & 0\\ 1 & -1\\ 1 & 1\\ -2 & 3
\end{pmatrix}
\qquad
b = \begin{pmatrix} 2\\ 1\\ 2\\ 1\\ 4\\ 6\end{pmatrix}
\qquad
c = \begin{pmatrix} 1 & 1\end{pmatrix}
\qquad
B = \{1,2\}
\]

\medskip
\hrule
\medskip

\textbf{Iterazione 1.} \quad $B=\{1,2\}$, $N=\{3,4,5,6\}$.
\[
A_B = \begin{pmatrix} -1 & -1\\ 0 & -1\end{pmatrix},\quad
b_B = \begin{pmatrix} 2\\ 1\end{pmatrix},
\qquad
\det A_B = 1,
\qquad
A_B^{-1} = \begin{pmatrix} -1 & 1\\ 0 & -1\end{pmatrix}
\]
\[
x = A_B^{-1} b_B = \begin{pmatrix} -1 & 1\\ 0 & -1\end{pmatrix}\begin{pmatrix}2\\1\end{pmatrix}
= \begin{pmatrix}-1\\-1\end{pmatrix}
\qquad c\,x = -2
\]
\[
y_B = c\,A_B^{-1} = \begin{pmatrix}1&1\end{pmatrix}\begin{pmatrix} -1 & 1\\ 0 & -1\end{pmatrix}
= \begin{pmatrix}-1 & 0\end{pmatrix}
\qquad\Longrightarrow\qquad y = (\,-1,\ 0,\ 0,\ 0,\ 0,\ 0\,)
\]
La prima componente è negativa: \textbf{non ottimo}. Esce $h = 1$ (prima posizione di $B$).
\[
\xi = -A_B^{-1} u_1 = -\begin{pmatrix}-1\\0\end{pmatrix} = \begin{pmatrix}1\\0\end{pmatrix}
\qquad\text{(prima colonna di } A_B^{-1}\text{, cambiata di segno)}
\]
\[
A_N \xi = \begin{pmatrix}1\\1\\1\\-2\end{pmatrix},
\qquad
A_N x = \begin{pmatrix}-1\\0\\-2\\-1\end{pmatrix},
\qquad
b_N = \begin{pmatrix}2\\1\\4\\6\end{pmatrix}
\]
Rapporti solo dove $A_i\xi > 0$ (quindi non $i=6$):
\[
\lambda_3 = \frac{2-(-1)}{1} = 3, \qquad
\lambda_4 = \frac{1-0}{1} = \mathbf{1}, \qquad
\lambda_5 = \frac{4-(-2)}{1} = 6
\]
Minimo: $k = 4$, $\lambda = 1$. Nuovo punto $x' = (-1,-1) + 1\cdot(1,0) = (0,-1)$.
Nuova base $B = \{2,4\}$.

\medskip
\hrule
\medskip

\textbf{Iterazione 2.} \quad $B=\{2,4\}$, $N=\{1,3,5,6\}$.
\[
A_B = \begin{pmatrix} 0 & -1\\ 1 & -1\end{pmatrix},\quad
b_B = \begin{pmatrix}1\\1\end{pmatrix},
\qquad \det A_B = 0\cdot(-1) - (-1)\cdot 1 = 1,
\qquad
A_B^{-1} = \begin{pmatrix} -1 & 1\\ -1 & 0\end{pmatrix}
\]
\[
x = \begin{pmatrix}-1&1\\-1&0\end{pmatrix}\begin{pmatrix}1\\1\end{pmatrix} = \begin{pmatrix}0\\-1\end{pmatrix}
\qquad c\,x = -1
\]
\[
y_B = \begin{pmatrix}1&1\end{pmatrix}\begin{pmatrix}-1&1\\-1&0\end{pmatrix} = \begin{pmatrix}-2 & 1\end{pmatrix}
\qquad\Longrightarrow\qquad y = (\,0,\ -2,\ 0,\ 1,\ 0,\ 0\,)
\]
Negativa la componente in prima posizione di $B$, cioè il vincolo $2$: esce $h = 2$.
\[
\xi = -A_B^{-1}u_1 = -\begin{pmatrix}-1\\-1\end{pmatrix} = \begin{pmatrix}1\\1\end{pmatrix}
\]
\[
A_N\xi = \begin{pmatrix}-2\\1\\2\\1\end{pmatrix},\qquad
A_N x = \begin{pmatrix}1\\0\\-1\\-3\end{pmatrix},\qquad
b_N = \begin{pmatrix}2\\2\\4\\6\end{pmatrix}
\qquad (N = \{1,3,5,6\})
\]
\[
\lambda_3 = \frac{2-0}{1} = \mathbf{2},\qquad
\lambda_5 = \frac{4-(-1)}{2} = 2.5,\qquad
\lambda_6 = \frac{6-(-3)}{1} = 9
\]
$k = 3$, $\lambda = 2$, $x' = (0,-1) + 2(1,1) = (2,1)$. Nuova base $B=\{3,4\}$.

\medskip
\hrule
\medskip

\textbf{Iterazione 3.} \quad $B=\{3,4\}$, $N=\{1,2,5,6\}$.
\[
A_B = \begin{pmatrix}1 & 0\\ 1 & -1\end{pmatrix},\quad
b_B = \begin{pmatrix}2\\1\end{pmatrix},
\qquad \det A_B = -1,
\qquad
A_B^{-1} = \frac{1}{-1}\begin{pmatrix}-1 & 0\\ -1 & 1\end{pmatrix}
= \begin{pmatrix}1 & 0\\ 1 & -1\end{pmatrix}
\]
\[
x = \begin{pmatrix}1&0\\1&-1\end{pmatrix}\begin{pmatrix}2\\1\end{pmatrix} = \begin{pmatrix}2\\1\end{pmatrix}
\qquad c\,x = 3
\]
\[
y_B = \begin{pmatrix}1&1\end{pmatrix}\begin{pmatrix}1&0\\1&-1\end{pmatrix} = \begin{pmatrix}2 & -1\end{pmatrix}
\qquad\Longrightarrow\qquad y = (\,0,\ 0,\ 2,\ -1,\ 0,\ 0\,)
\]
Negativa la \emph{seconda} componente, che corrisponde al vincolo $4$: esce $h = 4$.
\[
\xi = -A_B^{-1}u_2 = -\begin{pmatrix}0\\-1\end{pmatrix} = \begin{pmatrix}0\\1\end{pmatrix}
\]
\[
A_N\xi = \begin{pmatrix}-1\\-1\\1\\3\end{pmatrix},\qquad
A_N x = \begin{pmatrix}-3\\-1\\3\\-1\end{pmatrix},\qquad
b_N = \begin{pmatrix}2\\1\\4\\6\end{pmatrix}
\qquad (N=\{1,2,5,6\})
\]
\[
\lambda_5 = \frac{4-3}{1} = \mathbf{1},\qquad
\lambda_6 = \frac{6-(-1)}{3} = \tfrac{7}{3}
\]
$k = 5$, $\lambda = 1$, $x' = (2,1) + 1\cdot(0,1) = (2,2)$. Nuova base $B = \{3,5\}$.

\medskip
\hrule
\medskip

\textbf{Iterazione 4.} \quad $B=\{3,5\}$.
\[
A_B = \begin{pmatrix}1 & 0\\ 1 & 1\end{pmatrix},\quad
b_B = \begin{pmatrix}2\\4\end{pmatrix},
\qquad \det A_B = 1,
\qquad
A_B^{-1} = \begin{pmatrix}1 & 0\\ -1 & 1\end{pmatrix}
\]
\[
x = \begin{pmatrix}1&0\\-1&1\end{pmatrix}\begin{pmatrix}2\\4\end{pmatrix} = \begin{pmatrix}2\\2\end{pmatrix}
\qquad c\,x = 4
\]
\[
y_B = \begin{pmatrix}1&1\end{pmatrix}\begin{pmatrix}1&0\\-1&1\end{pmatrix} = \begin{pmatrix}0 & 1\end{pmatrix}
\ \ \ge 0
\qquad\Longrightarrow\qquad \textbf{ottimo.}
\]

\textbf{Risposta.} $x^* = (2,2)$, valore ottimo $z^* = 4$, certificato duale
$y^* = (0,0,0,0,1,0)$.

\textbf{Cammino percorso:} $(-1,-1) \to (0,-1) \to (2,1) \to (2,2)$; da riportare sul disegno
del punto (a).

\textbf{Osservazione da scrivere.} $y_3 = 0$ all'ottimo pur essendo il vincolo $3$ in base:
c'è quindi un \emph{intero spigolo} di soluzioni ottime, quello lungo il vincolo $5$
($x_1+x_2=4$), che è perpendicolare al gradiente $c=(1,1)$.
Infatti anche $(2,2)$ e ogni punto ammissibile con $x_1+x_2=4$ vale 4.
\end{es}

\begin{idea}
\textbf{Come presentarlo sul foglio.} Una tabella per iterazione, con:
$B$, $A_B$, $A_B^{-1}$, $x$, valore $c\,x$, $y_B$, vincolo uscente $h$, direzione $\xi$,
i rapporti $\lambda_i$, vincolo entrante $k$.
Poi ripassa sul disegno i vertici percorsi: è il cammino del simplesso.
\end{idea}

\subsection{Casi speciali da riconoscere}

\begin{center}
\begin{tabular}{@{}p{3.2cm} p{11.3cm}@{}}
\toprule
\textbf{Caso} & \textbf{Come si manifesta e cosa scrivere}\\
\midrule
Illimitato & Al passo dei rapporti, \emph{nessun} $i \in N$ ha $A_i\xi > 0$: la direzione $\xi$
non incontra alcun vincolo. Scrivi: «il problema è illimitato lungo la direzione $\xi = \dots$».\\
\addlinespace
Degenere & Più di $n$ vincoli attivi nello stesso vertice; risulta $\lambda_k = 0$: la base
cambia ma il punto no. Non è un errore, prosegui; usa Bland per non ciclare.\\
\addlinespace
Ottimi multipli & All'ottimo una componente di $y$ vale 0 su un vincolo di base: esiste
un intero spigolo di soluzioni ottime (vedi l'esempio sopra).\\
\addlinespace
Vuoto & Nessun punto soddisfa tutti i vincoli. All'esame non capita partendo da una base
ammissibile data.\\
\bottomrule
\end{tabular}
\end{center}

\subsection{Definizioni chieste a voce}

\begin{teo}{soluzione di base e soluzione di base ammissibile}
Dato il problema $\max c\,x$ con $Ax \le b$, $x \in \mathbb{R}^n$:

una \textbf{base} $B$ è un insieme di $n$ indici di vincoli le cui righe sono linearmente
indipendenti, cioè $\det A_B \ne 0$.

La \textbf{soluzione di base} associata a $B$ è l'unico punto $x_B = A_B^{-1} b_B$,
cioè il punto in cui quegli $n$ vincoli sono soddisfatti come uguaglianze (sono \emph{attivi}).

La soluzione di base è \textbf{ammissibile} se soddisfa anche tutti gli altri vincoli del
problema, cioè $A_N x_B \le b_N$. In tal caso essa coincide con un \textbf{vertice} del poliedro
delle soluzioni ammissibili.
\end{teo}

\begin{teo}{direzione ammissibile e direzione di crescita}
Sia $x$ una soluzione ammissibile e $I(x)$ l'insieme dei vincoli attivi in $x$.

Una direzione $\xi \ne 0$ è \textbf{ammissibile} in $x$ se esiste $\lambda > 0$ tale che
$x + \lambda \xi$ è ancora ammissibile. In forma $Ax\le b$ questo equivale a
$A_i \xi \le 0$ per ogni vincolo $i \in I(x)$ attivo in $x$.

La direzione è di \textbf{crescita} se $c\,\xi > 0$ (in un problema di massimo),
cioè se muoversi lungo $\xi$ aumenta il valore dell'obiettivo.

\textbf{Condizione di ottimalità.} $x$ è ottimo se e solo se in $x$ non esiste alcuna
direzione ammissibile che sia anche di crescita.
\end{teo}

\newpage
\subsection{Il duale}

Nella forma standard del corso:
\[
\text{Primale (P):} \ \max\ c\,x, \ \ A x \le b
\qquad\Longleftrightarrow\qquad
\text{Duale (D):} \ \min\ y\,b, \ \ y A = c, \ \ y \ge 0
\]
Il duale ha \textbf{una variabile per ogni vincolo del primale} ($y \in \mathbb{R}^m$, riga)
e \textbf{un vincolo per ogni variabile del primale} ($n$ equazioni).

\begin{es}{scrivere il duale --- compito 25/05/2026}
Primale (già normalizzato): $\max 2x_1 + x_2$ con
\[
\begin{aligned}
&(1)\ -x_1 \le 3 &&(3)\ x_1 - x_2 \le 1 &&(5)\ x_1 \le 4\\
&(2)\ -x_2 \le 2 &&(4)\ -x_1 + x_2 \le 1 &&(6)\ x_2 \le 4
\end{aligned}
\]
Duale: sei variabili $y_1,\dots,y_6 \ge 0$, due vincoli (uno per colonna di $A$):
\[
\min\ 3y_1 + 2y_2 + y_3 + y_4 + 4y_5 + 4y_6
\]
\[
\begin{aligned}
&-y_1 + y_3 - y_4 + y_5 = 2 &&\text{(colonna di } x_1)\\
&-y_2 - y_3 + y_4 + y_6 = 1 &&\text{(colonna di } x_2)\\
&y \ge 0
\end{aligned}
\]
\textbf{Come si costruisce meccanicamente:} il coefficiente di $y_i$ nel $j$-esimo vincolo duale
è $a_{ij}$, cioè si legge $A$ \emph{per colonne}; il termine noto del $j$-esimo vincolo duale è $c_j$;
i coefficienti dell'obiettivo duale sono i $b_i$.
\end{es}

\textbf{Tabella di conversione} (per problemi non nella forma standard):
\begin{center}
\begin{tabular}{@{}ll@{}}
\toprule
\textbf{Primale (massimo)} & \textbf{Duale (minimo)}\\
\midrule
vincolo $\le$              & variabile $\ge 0$\\
vincolo $\ge$              & variabile $\le 0$\\
vincolo $=$                & variabile libera\\
\addlinespace
variabile $\ge 0$          & vincolo $\ge c_j$\\
variabile libera           & vincolo $= c_j$\\
\bottomrule
\end{tabular}
\end{center}

\begin{teo}{dualità debole --- enunciato e dimostrazione}
\textbf{Enunciato.} Sia $x$ ammissibile per il primale $\max c\,x$, $Ax \le b$,
e sia $y \ge 0$ ammissibile per il duale $\min y\,b$, $yA = c$. Allora
\[ c\,x \ \le\ y\,b .\]
Cioè: \emph{ogni} soluzione ammissibile del duale fornisce un limite superiore al valore
di \emph{ogni} soluzione ammissibile del primale.

\textbf{Dimostrazione.}
\[ c\,x \;=\; (y A)\,x \;=\; y\,(A x) \;\le\; y\,b, \]
dove la prima uguaglianza usa l'ammissibilità duale $yA = c$, la seconda l'associatività,
e la disuguaglianza finale vale perché $Ax \le b$ e $y \ge 0$
(moltiplicare una disuguaglianza per un numero non negativo ne conserva il verso). $\square$

\textbf{Conseguenze (da citare sempre).}
\begin{itemize}
\item Se per una coppia $x, y$ ammissibili vale $c\,x = y\,b$, allora entrambe sono ottime.
      È esattamente ciò che certifica il simplesso quando si ferma con $y_B \ge 0$.
\item Se il primale è illimitato superiormente, il duale è vuoto (e viceversa).
\end{itemize}
\end{teo}

\textbf{Dualità forte.} Se uno dei due problemi ha ottimo finito, lo ha anche l'altro
e i due valori ottimi coincidono.

\textbf{Scarti complementari.} All'ottimo $y_i \,(b_i - A_i x) = 0$ per ogni $i$:
un vincolo non attivo ha variabile duale nulla, e una variabile duale positiva implica
vincolo attivo. È il motivo per cui nel simplesso $y$ vale zero fuori base.

\textbf{Interpretazione delle variabili duali (prezzi ombra).} $y_i^*$ misura di quanto varia
il valore ottimo se il termine noto $b_i$ cresce di una unità: è il «prezzo» che si sarebbe
disposti a pagare per una unità in più della risorsa $i$. Un vincolo non stringente ha
prezzo ombra nullo.

\newpage
\subsection{Branch-and-Bound}

\begin{ric}{Branch-and-Bound per un problema di massimo}
\begin{enumerate}
\item Risolvi il \textbf{rilassamento continuo} del nodo (togliendo l'interezza):
      dà un \textbf{bound superiore} $U_{\text{nodo}}$ per quel sottoproblema.
\item Se la soluzione trovata è già intera, è un candidato: aggiorna l'incumbent $L$
      (miglior soluzione intera finora) e \textbf{chiudi} il nodo.
\item Altrimenti scegli una variabile frazionaria $x_j^*$ e ramifica in due nodi figli:
      \[ x_j \le \lfloor x_j^* \rfloor \qquad\text{e}\qquad x_j \ge \lceil x_j^* \rceil \]
      (i due rami sono esaustivi e disgiunti: nessuna soluzione intera va persa).
\item \textbf{Chiudi} un nodo quando: è inammissibile, oppure dà soluzione intera,
      oppure $U_{\text{nodo}} \le L$ (\emph{dominanza}: non può contenere niente di meglio).
\item Quando l'albero è esaurito, $L$ è l'ottimo.
\end{enumerate}
Per un problema di minimo si scambiano i ruoli: bound inferiore $L_{\text{nodo}}$,
incumbent $U$, e si chiude quando $L_{\text{nodo}} \ge U$.
\end{ric}

\begin{es}{un albero completo con numeri piccoli}
\[ \max\ 5x_1 + 4x_2 \qquad 6x_1 + 4x_2 \le 24, \quad x_1 + 2x_2 \le 6, \quad x_1,x_2 \ge 0 \ \text{interi} \]

\textbf{Nodo 0 (radice).} Rilassamento continuo: l'ottimo è all'intersezione dei due vincoli.
Da $x_1 = 6-2x_2$ sostituito nel primo: $6(6-2x_2)+4x_2 = 24 \Rightarrow -8x_2 = -12 \Rightarrow x_2 = 1.5$,
$x_1 = 3$. Valore $U_0 = 15 + 6 = 21$. Non intera: si ramifica su $x_2$.

\textbf{Nodo 1: $x_2 \le 1$.} Con $x_2=1$: $6x_1 \le 20 \Rightarrow x_1 \le 10/3$;
$x_1 \le 4$. Ottimo continuo $x = (10/3,\,1)$, $U_1 = 50/3 + 4 \approx 20.67$. Frazionaria: si ramifica su $x_1$.
\begin{itemize}
\item \textbf{Nodo 1a: $x_1 \le 3$.} Ottimo $(3,1)$, valore $19$, \emph{intero} $\Rightarrow$
      incumbent $L = 19$. Chiuso.
\item \textbf{Nodo 1b: $x_1 \ge 4$.} $6\cdot 4 = 24$ forza $x_2 = 0$. Ottimo $(4,0)$, valore $20$,
      \emph{intero} e migliore $\Rightarrow$ incumbent $L = 20$. Chiuso.
\end{itemize}

\textbf{Nodo 2: $x_2 \ge 2$.} $x_1 + 2x_2 \le 6 \Rightarrow x_1 \le 2$;
ottimo continuo $(2,2)$ di valore $18$. Poiché $U_2 = 18 \le L = 20$, il nodo è
\textbf{chiuso per dominanza} (e comunque la soluzione è intera e peggiore).

\textbf{Conclusione.} Albero esaurito: ottimo intero $x^* = (4,0)$, $z^* = 20$.
Si noti che l'ottimo intero \emph{non} si ottiene arrotondando l'ottimo continuo $(3,\,1.5)$.
\end{es}

\newpage
%================================================================
\section{Esercizio 3 --- Reti di flusso}

\subsection{Glossario}

Rete $G = (N, A)$, con $u_{ij}$ capacità e $c_{ij}$ costo dell'arco $(i,j)$,
e bilancio $b_i$ al nodo.

\begin{att}{convenzione dei segni del bilancio}
Nei compiti di questo corso, nel problema di flusso di costo minimo il numero accanto al nodo è
il bilancio $b_i$ e vale: $b_i < 0$ per i nodi \textbf{sorgente} (offerta $|b_i|$),
$b_i > 0$ per i nodi \textbf{pozzo} (domanda $b_i$), $b_i = 0$ per i nodi di transito.
Nel compito del 27/07/2026: $A$ ha $-4$ e $B$ ha $-6$ (offrono 10 unità),
$D$ ha $3$ ed $E$ ha $7$ (ne chiedono 10). La somma dei bilanci deve fare zero.
\end{att}

\textbf{Flusso ammissibile.} Rispetta le capacità $0 \le x_{ij} \le u_{ij}$
e la conservazione ai nodi:
\[ \sum_{j :\, (i,j) \in A} x_{ij} \;-\; \sum_{j :\, (j,i) \in A} x_{ji} \;=\; -\,b_i \qquad \forall i \]
(uscite meno entrate al nodo $i$; con la convenzione dei segni sopra, un nodo sorgente
ha uscite nette positive).

\textbf{Pseudoflusso.} Rispetta le capacità $0\le x_{ij}\le u_{ij}$ ma \emph{non}
necessariamente la conservazione: i nodi possono avere \emph{eccessi} o \emph{difetti}.
Uno pseudoflusso è \textbf{ammissibile} quando tutti gli sbilanciamenti sono nulli,
cioè quando è un flusso.

\textbf{Taglio $s$-$t$.} Partizione $(S, T)$ dei nodi con $s \in S$ e $t \in T$.
La \textbf{capacità} conta solo gli archi in avanti:
\[ u(S,T) = \sum_{i \in S,\ j \in T} u_{ij} \]
Proprietà fondamentale: per ogni flusso $x$ e ogni taglio $(S,T)$,
$v(x) = \sum_{S\to T} x_{ij} - \sum_{T \to S} x_{ji} \le u(S,T)$.

\subsection{Grafo residuo $G_x$}

Per ogni arco $(i,j)$ con flusso corrente $x_{ij}$ si generano \emph{fino a due} archi residui:

\begin{center}
\begin{tabular}{@{}p{3.6cm} p{3.9cm} p{2.8cm} p{3.0cm}@{}}
\toprule
\textbf{Arco residuo} & \textbf{Esiste quando} & \textbf{Residuo} & \textbf{Costo}\\
\midrule
diretto \ $i \to j$ & $x_{ij} < u_{ij}$ & $u_{ij} - x_{ij}$ & $+\,c_{ij}$\\
\addlinespace
inverso \ $j \to i$ & $x_{ij} > 0$ & $x_{ij}$ & $-\,c_{ij}$\\
\bottomrule
\end{tabular}
\end{center}

\begin{idea}
\textbf{A che serve l'arco inverso.} Rappresenta la possibilità di \emph{disfare} flusso già
mandato. Percorrerlo significa «togli $\delta$ unità da $(i,j)$ e mandale altrove»: è per questo
che il flusso massimo trova l'ottimo anche se le prime scelte erano sbagliate.
Nel matching corrisponde a \emph{riassegnare} un accoppiamento già fatto.
\end{idea}

\newpage
\subsection{Flusso massimo con Edmonds--Karp}

\begin{ric}{Edmonds--Karp}
\begin{enumerate}
\item Parti dal flusso nullo $x = 0$.
\item Costruisci $G_x$ e cerca un cammino aumentante da $s$ a $t$ \textbf{con una visita BFS}
      (in ampiezza), cioè il cammino con il \emph{minor numero di archi}.
      A parità di distanza dichiara la regola usata per l'ordine dei nodi (alfabetico o numerico).
\item Calcola $\delta = \min$ dei residui lungo il cammino.
\item Aggiorna il flusso: $+\delta$ sugli archi percorsi in avanti, $-\delta$ su quelli percorsi
      all'indietro.
\item Ripeti finché $t$ non è più raggiungibile da $s$ in $G_x$.
\item \textbf{Taglio minimo.} $S$ = insieme dei nodi raggiungibili da $s$ nel $G_x$ finale,
      $T = N \setminus S$.
\end{enumerate}
\end{ric}

\begin{idea}
\textbf{Come si scrive una iterazione sul foglio.} Per ogni iterazione riporta:
(i) l'ordine di visita BFS con i predecessori, (ii) il cammino trovato, (iii) il $\delta$,
(iv) il nuovo valore del flusso, (v) il grafo aggiornato con le coppie
$u_{ij},\,x_{ij}$ su ogni arco.
\end{idea}

\begin{es}{Edmonds--Karp su una rete a 7 nodi}
Rete (le etichette sono le capacità):

\begin{center}
\begin{tikzpicture}[>=Latex,scale=1.15,transform shape,
  every node/.style={font=\small},
  nd/.style={circle,draw=navy,thick,minimum size=7mm,inner sep=1pt}]
\node[nd] (s) at (0,0)     {$s$};
\node[nd] (A) at (2.4,1.7) {$A$};
\node[nd] (B) at (2.4,-1.7){$B$};
\node[nd] (C) at (5.0,0)   {$C$};
\node[nd] (D) at (7.6,1.7) {$D$};
\node[nd] (E) at (7.6,-1.7){$E$};
\node[nd] (t) at (10.0,0)  {$t$};
\draw[->,thick] (s) -- node[above left]{4} (A);
\draw[->,thick] (s) -- node[below left]{6} (B);
\draw[->,thick] (A) -- node[above right,pos=0.4]{4} (C);
\draw[->,thick] (A) -- node[above]{6} (D);
\draw[->,thick] (B) -- node[left]{5} (A);
\draw[->,thick] (B) -- node[below right,pos=0.4]{4} (C);
\draw[->,thick] (B) to[bend right=22] node[below]{3} (E);
\draw[->,thick] (C) -- node[above right,pos=0.55]{4} (D);
\draw[->,thick] (C) -- node[below right,pos=0.55]{5} (E);
\draw[->,thick] (D) -- node[above right]{3} (t);
\draw[->,thick] (E) -- node[below right]{7} (t);
\end{tikzpicture}
\end{center}

Ordine dei nodi nella BFS: alfabetico ($A$ prima di $B$, ecc.).

\medskip
\textbf{Iterazione 1.} BFS da $s$: livello 1 $\{A, B\}$; livello 2 da $A$: $\{C, D\}$,
da $B$: $\{E\}$; livello 3 da $D$: $t$.
Cammino: $s \to A \to D \to t$. \quad
$\delta = \min(4,\,6,\,3) = 3$. \quad Flusso $= 3$.

\textbf{Iterazione 2.} Ora $D\to t$ è saturo. BFS: $s\to B \to E \to t$. \quad
$\delta = \min(6,\,3,\,7) = 3$. \quad Flusso $= 6$.

\textbf{Iterazione 3.} $B\to E$ è saturo. BFS: $s\to A \to C \to E \to t$.
Residui: $s\to A$ vale $4-3=1$, $A\to C$ vale 4, $C\to E$ vale 5, $E\to t$ vale $7-3=4$. \quad
$\delta = 1$. \quad Flusso $= 7$.

\textbf{Iterazione 4.} $s\to A$ è saturo. BFS: $s\to B \to C \to E \to t$.
Residui: $6-3=3$, $4$, $5-1=4$, $7-4=3$. \quad $\delta = 3$. \quad Flusso $= 10$.

\textbf{Iterazione 5.} In $G_x$ da $s$ non si raggiunge più nulla: $s\to A$ ($4/4$) e
$s\to B$ ($6/6$) sono entrambi saturi. \textbf{Stop.}

\medskip
\textbf{Soluzione.} Flusso massimo $= 10$, con
\[
x_{sA}=4,\quad x_{sB}=6,\quad x_{AC}=1,\quad x_{AD}=3,\quad x_{BA}=0,\quad x_{BC}=3,
\]
\[
x_{BE}=3,\quad x_{CD}=0,\quad x_{CE}=4,\quad x_{Dt}=3,\quad x_{Et}=7.
\]

\textbf{Taglio minimo.} $S = \{s\}$, $T = \{A,B,C,D,E,t\}$.
Verifica: archi $S \to T$ sono $s\to A$ e $s\to B$, entrambi saturi;
archi $T \to S$: nessuno. Capacità $u(S,T) = 4 + 6 = 10 = v(x)$. \checkmark
\end{es}

\begin{idea}
\textbf{Verifica finale del taglio (tre controlli, sempre da scrivere).}
(1) gli archi da $S$ a $T$ sono tutti saturi; (2) gli archi da $T$ a $S$ sono tutti vuoti;
(3) $u(S,T)$ è pari al valore del flusso trovato.
Se i tre controlli tornano, la soluzione è certamente ottima.
\end{idea}

\textbf{Complessità e differenza con Ford--Fulkerson.}
Ford--Fulkerson generico non specifica come scegliere il cammino aumentante e costa
$O(m \cdot U)$: dipende dai \emph{valori} delle capacità e su istanze costruite ad arte
degenera. Edmonds--Karp è la specializzazione che sceglie sempre il cammino aumentante con
il minor numero di archi, esplorando $G_x$ \textbf{in ampiezza (BFS)}, e costa $O(n\,m^2)$,
\emph{indipendentemente} dai valori delle capacità. È questa la frase da scrivere quando
la domanda è «in cosa differisce Edmonds--Karp dalle altre implementazioni».

\newpage
\begin{teo}{teorema max-flow-min-cut}
\textbf{Enunciato.} Il valore del flusso massimo da $s$ a $t$ è uguale alla capacità
del taglio $s$-$t$ di capacità minima.

\textbf{Parte 1: ogni flusso $\le$ ogni taglio.}
Per ogni flusso ammissibile $x$ e ogni taglio $(S,T)$ vale
\[
v(x) \;=\; \sum_{i\in S,\,j\in T} x_{ij} \;-\; \sum_{i\in T,\,j\in S} x_{ij}
      \;\le\; \sum_{i\in S,\,j\in T} x_{ij}
      \;\le\; \sum_{i\in S,\,j\in T} u_{ij} \;=\; u(S,T),
\]
dove la prima disuguaglianza vale perché il flusso all'indietro è $\ge 0$ e la seconda
perché $x_{ij}\le u_{ij}$.
Quindi $\max_x v(x) \le \min_{(S,T)} u(S,T)$.

\textbf{Parte 2: esistono un flusso e un taglio che si toccano.}
All'arresto di Ford--Fulkerson (nessun cammino aumentante) sia
$S$ l'insieme dei nodi raggiungibili da $s$ in $G_x$; per definizione di arresto $t \notin S$,
quindi $(S,T)$ con $T = N\setminus S$ è un taglio $s$-$t$.
Ogni arco $(i,j)$ con $i\in S$, $j\in T$ è \emph{saturo} ($x_{ij} = u_{ij}$): altrimenti
esisterebbe l'arco residuo diretto $i\to j$ e $j$ sarebbe raggiungibile, contro $j\in T$.
Ogni arco $(j,i)$ con $j\in T$, $i\in S$ è \emph{vuoto} ($x_{ji} = 0$): altrimenti esisterebbe
l'arco residuo inverso $i\to j$ e di nuovo $j$ sarebbe raggiungibile.
Sostituendo nella formula della parte 1:
\[ v(x) = \sum_{S\to T} u_{ij} - 0 = u(S,T). \]

\textbf{Conclusione.} Per la parte 1 nessun flusso può superare $u(S,T)$ e nessun taglio può
scendere sotto $v(x)$; poiché i due valori coincidono, $x$ è un flusso massimo e $(S,T)$ un
taglio minimo. $\square$
\end{teo}

\subsection{Flusso di costo minimo}

Due algoritmi, entrambi possibili all'esame. Il primo è quello effettivamente chiesto
nel compito del 27/07/2026.

\begin{ric}{cancellazione dei cicli --- algoritmo primale}
\begin{enumerate}
\item Trova un flusso \textbf{ammissibile} qualsiasi, che rispetti i bilanci $b_i$.
      Il modo standard: aggiungi una sorgente fittizia $s$ collegata a ogni nodo di offerta
      (capacità $=$ offerta) e un pozzo fittizio $t$ collegato da ogni nodo di domanda
      (capacità $=$ domanda), poi calcola il \textbf{flusso massimo}.
      Se il flusso massimo satura tutti gli archi fittizi, esiste un flusso ammissibile
      ed è quello trovato; altrimenti il problema è inammissibile.
\item Costruisci $G_x$ con i costi: $+c_{ij}$ sugli archi diretti, $-c_{ij}$ sugli inversi.
\item Cerca un \textbf{ciclo orientato di costo negativo} in $G_x$, per ispezione o con
      Bellman--Ford.
\item Calcola $\delta = \min$ dei residui sul ciclo e spingi $\delta$ lungo il ciclo:
      il costo totale cala di $\delta \cdot |c(\text{ciclo})|$.
\item Ripeti finché non restano cicli negativi: allora il flusso è ottimo.
\end{enumerate}
\end{ric}

\begin{idea}
\textbf{Come si calcola il costo di un ciclo.} Si sommano i costi degli archi \emph{residui}
percorsi: $+c_{ij}$ se percorri l'arco nel verso originale, $-c_{ij}$ se lo percorri
all'indietro (stai disfacendo flusso, quindi recuperi il costo).
\end{idea}

\newpage
\begin{es}{cancellazione dei cicli --- compito del 27/07/2026}

\textbf{Rete di partenza.} Accanto a ogni arco la coppia $(\text{capacità } u_{ij},\, \text{costo } c_{ij})$; accanto ai nodi il bilancio $b_i$.

\begin{center}
\begin{tikzpicture}[>=Latex,scale=1.15,transform shape,
  every node/.style={font=\small},
  nd/.style={circle,draw=navy,thick,minimum size=8mm,inner sep=1pt}]
\node[nd] (A) at (0,2.2) {$A$};
\node[nd] (B) at (0,-2.2) {$B$};
\node[nd] (C) at (3.2,0) {$C$};
\node[nd] (D) at (6.4,2.2) {$D$};
\node[nd] (E) at (6.4,-2.2) {$E$};
\node[left=2mm of A,text=crimson] {$-4$};
\node[left=2mm of B,text=crimson] {$-6$};
\node[right=2mm of D,text=teal] {$3$};
\node[right=2mm of E,text=teal] {$7$};
\draw[->,thick] (A) -- node[above]{$6,\,6$} (D);
\draw[->,thick] (A) -- node[above right,pos=0.45]{$4,\,1$} (C);
\draw[->,thick] (C) -- node[above right,pos=0.55]{$4,\,4$} (D);
\draw[->,thick] (B) -- node[left]{$5,\,1$} (A);
\draw[->,thick] (B) -- node[below right,pos=0.45]{$4,\,1$} (C);
\draw[->,thick] (C) -- node[below right,pos=0.55]{$5,\,2$} (E);
\draw[->,thick] (B) -- node[below]{$3,\,6$} (E);
\draw[->,thick] (D) -- node[right]{$3,\,-3$} (E);
\end{tikzpicture}
\end{center}

% --- SCHEDA OPERATIVA STEP-BY-STEP ---
\begin{tcolorbox}[colback=lightbg,colframe=navy,arc=2pt,boxrule=0.8pt,left=6pt,right=6pt,top=5pt,bottom=5pt]
\textbf{Metodo pratico d'esame: come tenere traccia delle variabili senza confondersi}
\begin{enumerate}[leftmargin=1.2em,itemsep=0.2em]
\item \textbf{Sul grafo originale $G$}: tieni traccia solo di $[x_{ij} / u_{ij}]$ (flusso inviato / capacità).
\item \textbf{Sul grafo residuo $G_x$}: si disegnano solo gli archi che hanno residuo $r > 0$:
  \begin{itemize}[leftmargin=1.1em,itemsep=0.1em]
    \item \textbf{Arco Diretto} ($i \to j$, verso originale): esiste se $x_{ij} < u_{ij}$.\\
          Residuo disponibile: $r_{ij} = u_{ij} - x_{ij}$ \quad | \quad Costo residuo: $+c_{ij}$.
    \item \textbf{Arco Inverso} ($j \to i$, verso opposto): esiste solo se c'è flusso da disfare, cioè $x_{ij} > 0$.\\
          Residuo disponibile: $r_{ji} = x_{ij}$ \quad | \quad Costo residuo: $-c_{ij}$ \emph{(recuperi il costo già pagato)}.
    \item Se $0 < x_{ij} < u_{ij}$, esistono \textbf{entrambi} i versi residui ($\leftrightarrow$).
  \end{itemize}
\item \textbf{Identificazione ciclo $\Gamma$}: cerca un ciclo orientato chiuso in $G_x$ con $\sum c^{\text{res}} < 0$.
\item \textbf{Quantità di incremento $\delta$}: collo di bottiglia dei residui lungo il ciclo, $\delta = \min \{r_{ij}\}$.
\item \textbf{Aggiornamento del flusso}: percorrendo il ciclo $\Gamma$, per ogni arco attraversato:
  \[
  \begin{cases}
  \text{se percorso in verso \textbf{diretto}}: & x_{ij} \leftarrow x_{ij} + \delta\\
  \text{se percorso in verso \textbf{inverso}}: & x_{ij} \leftarrow x_{ij} - \delta
  \end{cases}
  \qquad\Longrightarrow\quad \Delta \text{Costo} = \delta \cdot c(\Gamma)
  \]
\end{enumerate}
\end{tcolorbox}

\textbf{Passo 0 --- Flusso ammissibile di partenza.}
Aggiungendo una sorgente ausiliaria $s$ verso $A, B$ e un pozzo $t$ da $D, E$, l'algoritmo di Edmonds--Karp determina il flusso iniziale ammissibile di valore 10:
\[
x_{AC}=1,\quad x_{AD}=3,\quad x_{BA}=0,\quad x_{BC}=3,\quad x_{BE}=3,\quad x_{CD}=0,\quad x_{CE}=4,\quad x_{DE}=0
\]
Costo iniziale: $z_0 = (1\cdot 1) + (3\cdot 6) + (0\cdot 1) + (3\cdot 1) + (3\cdot 6) + (0\cdot 4) + (4\cdot 2) + (0\cdot -3) = \mathbf{48}$.

\medskip
\hrule
\medskip

% ==================== ITERAZIONE 1 ====================
\textbf{Iterazione 1.}

\textbf{1. Tracciamento del ciclo negativo $\Gamma_1 = D \to E \to C \to D$:}
\begin{center}
\small
\begin{tabular}{@{}ccccc@{}}
\toprule
\textbf{Arco residuo} & \textbf{Tipo} & \textbf{Formula residuo $r$} & \textbf{Residuo num.} & \textbf{Costo residuo $c^{\text{res}}$}\\
\midrule
$D \to E$ & Diretto & $u_{DE} - x_{DE} = 3 - 0$ & $3$ & $+(-3) = -3$\\
$E \to C$ & Inverso & $x_{CE}$ & $4$ & $-(+2) = -2$\\
$C \to D$ & Diretto & $u_{CD} - x_{CD} = 4 - 0$ & $4$ & $+(+4) = +4$\\
\bottomrule
\end{tabular}
\end{center}
\begin{itemize}[itemsep=0.15em]
\item \textbf{Costo del ciclo}: $c(\Gamma_1) = (-3) + (-2) + (+4) = \mathbf{-1} < 0$ (ciclo migliorante).
\item \textbf{Capacità di incremento}: $\delta = \min(3, 4, 4) = \mathbf{3}$.
\item \textbf{Aggiornamento flussi}:
  $x_{DE} \leftarrow 0 + 3 = \mathbf{3}$ (saturo), \quad
  $x_{CE} \leftarrow 4 - 3 = \mathbf{1}$, \quad
  $x_{CD} \leftarrow 0 + 3 = \mathbf{3}$.
\item \textbf{Nuovo costo}: $z_1 = z_0 + \delta \cdot c(\Gamma_1) = 48 + 3(-1) = \mathbf{45}$.
\end{itemize}

\begin{center}
\begin{minipage}[t]{0.48\textwidth}\centering
\begin{tikzpicture}[>=Latex,scale=0.75,transform shape]
\mcfnodes
\ra{<->}{A}{D}{+6}
\ra{<->}{A}{C}{+1}
\ra{->}{C}{D}{+4}
\ra{->}{B}{A}{+1}
\ra{<->}{B}{C}{+1}
\ra{<->}{C}{E}{+2}
\ra{<-}{B}{E}{+6}
\ra{->}{D}{E}{-3}
\cy{D}{E}{-3}
\cy{E}{C}{-2}
\cy{C}{D}{+4}
\end{tikzpicture}
\par\vspace{1mm}{\scriptsize $G_x$ iniziale: evidenziato $\Gamma_1$ ($c=-1$, $\delta=3$)}
\end{minipage}\hfill
\begin{minipage}[t]{0.48\textwidth}\centering
\begin{tikzpicture}[>=Latex,scale=0.75,transform shape]
\mcfnodes
\fa{A}{D}{6,\,6}{3}
\fa{A}{C}{4,\,1}{1}
\fa{C}{D}{4,\,4}{3}
\fa{B}{A}{5,\,1}{0}
\fa{B}{C}{4,\,1}{3}
\fa{C}{E}{5,\,2}{1}
\fa{B}{E}{3,\,6}{3}
\fa{D}{E}{3,\,-3}{3}
\end{tikzpicture}
\par\vspace{1mm}{\scriptsize $G$ dopo iterazione 1 (costo $45$)}
\end{minipage}
\end{center}

\medskip
\hrule
\medskip

% ==================== ITERAZIONE 2 ====================
\textbf{Iterazione 2.}

\textbf{1. Tracciamento del ciclo negativo $\Gamma_2 = C \to E \to B \to C$:}
\begin{center}
\small
\begin{tabular}{@{}ccccc@{}}
\toprule
\textbf{Arco residuo} & \textbf{Tipo} & \textbf{Formula residuo $r$} & \textbf{Residuo num.} & \textbf{Costo residuo $c^{\text{res}}$}\\
\midrule
$C \to E$ & Diretto & $u_{CE} - x_{CE} = 5 - 1$ & $4$ & $+2$\\
$E \to B$ & Inverso & $x_{BE}$ & $3$ & $-6$\\
$B \to C$ & Diretto & $u_{BC} - x_{BC} = 4 - 3$ & $1$ & $+1$\\
\bottomrule
\end{tabular}
\end{center}
\begin{itemize}[itemsep=0.15em]
\item \textbf{Costo del ciclo}: $c(\Gamma_2) = 2 - 6 + 1 = \mathbf{-3} < 0$.
\item \textbf{Capacità di incremento}: $\delta = \min(4, 3, 1) = \mathbf{1}$.
\item \textbf{Aggiornamento flussi}:
  $x_{CE} \leftarrow 1 + 1 = \mathbf{2}$, \quad
  $x_{BE} \leftarrow 3 - 1 = \mathbf{2}$, \quad
  $x_{BC} \leftarrow 3 + 1 = \mathbf{4}$ (saturo).
\item \textbf{Nuovo costo}: $z_2 = 45 + 1(-3) = \mathbf{42}$.
\end{itemize}

\begin{center}
\begin{minipage}[t]{0.48\textwidth}\centering
\begin{tikzpicture}[>=Latex,scale=0.75,transform shape]
\mcfnodes
\ra{<->}{A}{D}{+6}
\ra{<->}{A}{C}{+1}
\ra{<->}{C}{D}{+4}
\ra{->}{B}{A}{+1}
\ra{<->}{B}{C}{+1}
\ra{<->}{C}{E}{+2}
\ra{<-}{B}{E}{+6}
\ra{<-}{D}{E}{-3}
\cy{C}{E}{+2}
\cy{E}{B}{-6}
\cy{B}{C}{+1}
\end{tikzpicture}
\par\vspace{1mm}{\scriptsize $G_x$ post it. 1: evidenziato $\Gamma_2$ ($c=-3$, $\delta=1$)}
\end{minipage}\hfill
\begin{minipage}[t]{0.48\textwidth}\centering
\begin{tikzpicture}[>=Latex,scale=0.75,transform shape]
\mcfnodes
\fa{A}{D}{6,\,6}{3}
\fa{A}{C}{4,\,1}{1}
\fa{C}{D}{4,\,4}{3}
\fa{B}{A}{5,\,1}{0}
\fa{B}{C}{4,\,1}{4}
\fa{C}{E}{5,\,2}{2}
\fa{B}{E}{3,\,6}{2}
\fa{D}{E}{3,\,-3}{3}
\end{tikzpicture}
\par\vspace{1mm}{\scriptsize $G$ dopo iterazione 2 (costo $42$)}
\end{minipage}
\end{center}

\medskip
\hrule
\medskip

% ==================== ITERAZIONE 3 ====================
\textbf{Iterazione 3.}

\textbf{1. Tracciamento del ciclo negativo $\Gamma_3 = C \to E \to B \to A \to C$:}
\begin{center}
\small
\begin{tabular}{@{}ccccc@{}}
\toprule
\textbf{Arco residuo} & \textbf{Tipo} & \textbf{Formula residuo $r$} & \textbf{Residuo num.} & \textbf{Costo residuo $c^{\text{res}}$}\\
\midrule
$C \to E$ & Diretto & $u_{CE} - x_{CE} = 5 - 2$ & $3$ & $+2$\\
$E \to B$ & Inverso & $x_{BE}$ & $2$ & $-6$\\
$B \to A$ & Diretto & $u_{BA} - x_{BA} = 5 - 0$ & $5$ & $+1$\\
$A \to C$ & Diretto & $u_{AC} - x_{AC} = 4 - 1$ & $3$ & $+1$\\
\bottomrule
\end{tabular}
\end{center}
\begin{itemize}[itemsep=0.15em]
\item \textbf{Costo del ciclo}: $c(\Gamma_3) = 2 - 6 + 1 + 1 = \mathbf{-2} < 0$.
\item \textbf{Capacità di incremento}: $\delta = \min(3, 2, 5, 3) = \mathbf{2}$.
\item \textbf{Aggiornamento flussi}:
  $x_{CE} \leftarrow 2 + 2 = \mathbf{4}$, \quad
  $x_{BE} \leftarrow 2 - 2 = \mathbf{0}$ (vuoto), \quad
  $x_{BA} \leftarrow 0 + 2 = \mathbf{2}$, \quad
  $x_{AC} \leftarrow 1 + 2 = \mathbf{3}$.
\item \textbf{Nuovo costo}: $z_3 = 42 + 2(-2) = \mathbf{38}$.
\end{itemize}

\begin{center}
\begin{minipage}[t]{0.48\textwidth}\centering
\begin{tikzpicture}[>=Latex,scale=0.75,transform shape]
\mcfnodes
\ra{<->}{A}{D}{+6}
\ra{<->}{A}{C}{+1}
\ra{<->}{C}{D}{+4}
\ra{->}{B}{A}{+1}
\ra{<-}{B}{C}{+1}
\ra{<->}{C}{E}{+2}
\ra{<->}{B}{E}{+6}
\ra{<-}{D}{E}{-3}
\cy{C}{E}{+2}
\cy{E}{B}{-6}
\cy{B}{A}{+1}
\cy{A}{C}{+1}
\end{tikzpicture}
\par\vspace{1mm}{\scriptsize $G_x$ post it. 2: evidenziato $\Gamma_3$ ($c=-2$, $\delta=2$)}
\end{minipage}\hfill
\begin{minipage}[t]{0.48\textwidth}\centering
\begin{tikzpicture}[>=Latex,scale=0.75,transform shape]
\mcfnodes
\fa{A}{D}{6,\,6}{3}
\fa{A}{C}{4,\,1}{3}
\fa{C}{D}{4,\,4}{3}
\fa{B}{A}{5,\,1}{2}
\fa{B}{C}{4,\,1}{4}
\fa{C}{E}{5,\,2}{4}
\fa{B}{E}{3,\,6}{0}
\fa{D}{E}{3,\,-3}{3}
\end{tikzpicture}
\par\vspace{1mm}{\scriptsize $G$ dopo iterazione 3 (costo $38$)}
\end{minipage}
\end{center}

\medskip
\hrule
\medskip

% ==================== ITERAZIONE 4 ====================
\textbf{Iterazione 4.}

\textbf{1. Tracciamento del ciclo negativo $\Gamma_4 = C \to D \to A \to C$:}
\begin{center}
\small
\begin{tabular}{@{}ccccc@{}}
\toprule
\textbf{Arco residuo} & \textbf{Tipo} & \textbf{Formula residuo $r$} & \textbf{Residuo num.} & \textbf{Costo residuo $c^{\text{res}}$}\\
\midrule
$C \to D$ & Diretto & $u_{CD} - x_{CD} = 4 - 3$ & $1$ & $+4$\\
$D \to A$ & Inverso & $x_{AD}$ & $3$ & $-6$\\
$A \to C$ & Diretto & $u_{AC} - x_{AC} = 4 - 3$ & $1$ & $+1$\\
\bottomrule
\end{tabular}
\end{center}
\begin{itemize}[itemsep=0.15em]
\item \textbf{Costo del ciclo}: $c(\Gamma_4) = 4 - 6 + 1 = \mathbf{-1} < 0$.
\item \textbf{Capacità di incremento}: $\delta = \min(1, 3, 1) = \mathbf{1}$.
\item \textbf{Aggiornamento flussi}:
  $x_{CD} \leftarrow 3 + 1 = \mathbf{4}$ (saturo), \quad
  $x_{AD} \leftarrow 3 - 1 = \mathbf{2}$, \quad
  $x_{AC} \leftarrow 3 + 1 = \mathbf{4}$ (saturo).
\item \textbf{Nuovo costo}: $z_4 = 38 + 1(-1) = \mathbf{37}$.
\end{itemize}

\begin{center}
\begin{minipage}[t]{0.48\textwidth}\centering
\begin{tikzpicture}[>=Latex,scale=0.75,transform shape]
\mcfnodes
\ra{<->}{A}{D}{+6}
\ra{<->}{A}{C}{+1}
\ra{<->}{C}{D}{+4}
\ra{<->}{B}{A}{+1}
\ra{<-}{B}{C}{+1}
\ra{<->}{C}{E}{+2}
\ra{->}{B}{E}{+6}
\ra{<-}{D}{E}{-3}
\cy{C}{D}{+4}
\cy{D}{A}{-6}
\cy{A}{C}{+1}
\end{tikzpicture}
\par\vspace{1mm}{\scriptsize $G_x$ post it. 3: evidenziato $\Gamma_4$ ($c=-1$, $\delta=1$)}
\end{minipage}\hfill
\begin{minipage}[t]{0.48\textwidth}\centering
\begin{tikzpicture}[>=Latex,scale=0.75,transform shape]
\mcfnodes
\fa{A}{D}{6,\,6}{2}
\fa{A}{C}{4,\,1}{4}
\fa{C}{D}{4,\,4}{4}
\fa{B}{A}{5,\,1}{2}
\fa{B}{C}{4,\,1}{4}
\fa{C}{E}{5,\,2}{4}
\fa{B}{E}{3,\,6}{0}
\fa{D}{E}{3,\,-3}{3}
\end{tikzpicture}
\par\vspace{1mm}{\scriptsize $G$ dopo iterazione 4 (costo $37$)}
\end{minipage}
\end{center}

\medskip
\hrule
\medskip

% ==================== ITERAZIONE 5 (OTTIMO) ====================
\textbf{Iterazione 5 (Test di arresto ed ottimalità).}

Nel nuovo grafo residuo $G_x$, gli unici cicli orientati rimasti sono:
\begin{itemize}[itemsep=0.1em]
\item $A \to D \to E \to C \to A$ con costo: $(+6) + (-3) + (-2) + (-1) = 0 \ge 0$.
\item $B \to A \to D \to E \to B$ con costo: $(+1) + (+6) + (-3) + (+6) = +10 \ge 0$.
\end{itemize}
Non esiste alcun ciclo con costo strettamente negativo in $G_x$: \textbf{la soluzione corrente è ottima}.

\begin{center}
\begin{minipage}[t]{0.48\textwidth}\centering
\begin{tikzpicture}[>=Latex,scale=0.75,transform shape]
\mcfnodes
\ra{<->}{A}{D}{+6}
\ra{<-}{A}{C}{+1}
\ra{<-}{C}{D}{+4}
\ra{<->}{B}{A}{+1}
\ra{<-}{B}{C}{+1}
\ra{<->}{C}{E}{+2}
\ra{->}{B}{E}{+6}
\ra{<-}{D}{E}{-3}
\end{tikzpicture}
\par\vspace{1mm}{\scriptsize $G_x$ finale: privo di cicli orientati negativi}
\end{minipage}\hfill
\begin{minipage}[t]{0.48\textwidth}\centering
\begin{tikzpicture}[>=Latex,scale=0.75,transform shape]
\mcfnodes
\fa{A}{D}{6,\,6}{2}
\fa{A}{C}{4,\,1}{4}
\fa{C}{D}{4,\,4}{4}
\fa{B}{A}{5,\,1}{2}
\fa{B}{C}{4,\,1}{4}
\fa{C}{E}{5,\,2}{4}
\fa{B}{E}{3,\,6}{0}
\fa{D}{E}{3,\,-3}{3}
\end{tikzpicture}
\par\vspace{1mm}{\scriptsize Flusso ottimo $x^*$, costo minimo $\mathbf{37}$}
\end{minipage}
\end{center}

\medskip
\hrule
\medskip

\textbf{Soluzione ottima finale.}
\[
x_{AC}^* = 4,\quad x_{AD}^* = 2,\quad x_{BA}^* = 2,\quad x_{BC}^* = 4,\quad x_{BE}^* = 0,
\quad x_{CD}^* = 4,\quad x_{CE}^* = 4,\quad x_{DE}^* = 3
\]
\textbf{Costo totale minimo}:
\[
z^* = (4\cdot 1) + (2\cdot 6) + (2\cdot 1) + (4\cdot 1) + (0\cdot 6) + (4\cdot 4) + (4\cdot 2) + (3\cdot -3) = \mathbf{37}
\]
\textbf{Verifica finale dei bilanci ai nodi} (uscite $-$ entrate $= -b_i$):
\begin{itemize}[itemsep=0.15em]
\item Nodo $A$: $(x_{AC} + x_{AD}) - x_{BA} = (4 + 2) - 2 = 4 = -(-4)$ \quad \checkmark
\item Nodo $B$: $(x_{BA} + x_{BC} + x_{BE}) - 0 = (2 + 4 + 0) = 6 = -(-6)$ \quad \checkmark
\item Nodo $C$: $(x_{CD} + x_{CE}) - (x_{AC} + x_{BC}) = (4 + 4) - (4 + 4) = 0 = -(0)$ \quad \checkmark
\item Nodo $D$: $x_{DE} - (x_{AD} + x_{CD}) = 3 - (2 + 4) = -3 = -(3)$ \quad \checkmark
\item Nodo $E$: $0 - (x_{BE} + x_{CE} + x_{DE}) = 0 - (0 + 4 + 3) = -7 = -(7)$ \quad \checkmark
\end{itemize}
\end{es}

\begin{ric}{cammini minimi successivi --- algoritmo duale}
\begin{enumerate}
\item Parti da uno \textbf{pseudoflusso minimale}, cioè già ottimo nei costi
      ma non necessariamente ammissibile.
      Se tutti i costi sono non negativi va bene $x = 0$: nessun ciclo negativo,
      ma la conservazione è violata ovunque.
\item Finché esistono nodi sbilanciati, scegli $s$ con eccesso positivo e $t$ con difetto.
\item Cerca in $G_x$ il \textbf{cammino di costo minimo} da $s$ a $t$ e poni
      \[ \delta = \min\big(\text{eccesso di } s,\ |\text{difetto di } t|,\ \text{residui del cammino}\big) \]
\item Spingi $\delta$ lungo il cammino e aggiorna eccessi e difetti.
      Termina quando tutti i bilanci sono soddisfatti: quel flusso è ottimo.
\end{enumerate}
\end{ric}

\begin{es}{cammini minimi successivi --- stessa rete del compito del 27/07/2026}
Stessa rete dell'esempio di cancellazione dei cicli (capacità, costi e bilanci identici):
si arriva allo stesso ottimo $z^*=37$, ma partendo da uno pseudoflusso invece che da un flusso ammissibile.

\textbf{Rete di partenza.} Coppia $(u_{ij},\,c_{ij})$ sugli archi, bilancio $b_i$ accanto ai nodi.
\begin{center}
\begin{tikzpicture}[>=Latex,scale=1.15,transform shape,
  every node/.style={font=\small},
  nd/.style={circle,draw=navy,thick,minimum size=8mm,inner sep=1pt}]
\node[nd] (A) at (0,2.2) {$A$};
\node[nd] (B) at (0,-2.2) {$B$};
\node[nd] (C) at (3.2,0) {$C$};
\node[nd] (D) at (6.4,2.2) {$D$};
\node[nd] (E) at (6.4,-2.2) {$E$};
\node[left=2mm of A,text=crimson] {$-4$};
\node[left=2mm of B,text=crimson] {$-6$};
\node[right=2mm of D,text=teal] {$3$};
\node[right=2mm of E,text=teal] {$7$};
\draw[->,thick] (A) -- node[above]{$6,\,6$} (D);
\draw[->,thick] (A) -- node[above right,pos=0.45]{$4,\,1$} (C);
\draw[->,thick] (C) -- node[above right,pos=0.55]{$4,\,4$} (D);
\draw[->,thick] (B) -- node[left]{$5,\,1$} (A);
\draw[->,thick] (B) -- node[below right,pos=0.45]{$4,\,1$} (C);
\draw[->,thick] (C) -- node[below right,pos=0.55]{$5,\,2$} (E);
\draw[->,thick] (B) -- node[below]{$3,\,6$} (E);
\draw[->,thick] (D) -- node[right]{$3,\,-3$} (E);
\end{tikzpicture}
\end{center}

\begin{tcolorbox}[colback=lightbg,colframe=navy,arc=2pt,boxrule=0.8pt,left=6pt,right=6pt,top=5pt,bottom=5pt]
\textbf{Metodo pratico d'esame: come tenere traccia degli sbilanciamenti}
\begin{enumerate}[leftmargin=1.2em,itemsep=0.2em]
\item \textbf{Pseudoflusso di partenza}: $x_{ij}=0$ sugli archi con $c_{ij}\ge 0$ e
      $x_{ij}=u_{ij}$ sugli archi con $c_{ij}<0$ (\emph{saturare gli archi negativi} è
      indispensabile: altrimenti in $G_x$ resta un arco di costo negativo e lo pseudoflusso non è minimale).
\item \textbf{Sbilanciamento di ogni nodo}: posto $s_i=-b_i$ (offerta del nodo),
      \[ e_i \;=\; s_i \;+\; \text{(flusso entrante)} \;-\; \text{(flusso uscente)} \]
      $e_i>0$ = \textbf{eccesso} (sorgente da svuotare), $e_i<0$ = \textbf{difetto} (pozzo da riempire).
\item \textbf{Scelta della coppia}: prendi un $s$ con $e_s>0$ e un $t$ con $e_t<0$ e cerca in $G_x$
      il cammino di \emph{costo minimo} (non il più corto in numero di archi!).
\item \textbf{Passo}: $\delta = \min\big(e_s,\ |e_t|,\ \min r_{ij}\text{ sul cammino}\big)$; spingi $\delta$,
      aggiorna $x$ (\,$+\delta$ sugli archi diretti, $-\delta$ su quelli inversi\,), aggiorna $e_s$ ed $e_t$.
\item \textbf{Terminazione}: quando tutti gli $e_i$ sono nulli lo pseudoflusso è diventato un flusso,
      ed è già ottimo (ogni passo conserva la minimalità).
\end{enumerate}
\end{tcolorbox}

\textbf{Passo 0 --- Pseudoflusso minimale iniziale.}
L'unico arco di costo negativo è $D\to E$ ($c=-3$): lo saturo. Tutti gli altri restano a zero.
\[
x_{DE}=3,\qquad x_{ij}=0 \text{ per ogni altro arco}
\qquad\Longrightarrow\qquad z_0 = 3\cdot(-3) = \mathbf{-9}
\]
Sbilanciamenti (offerte: $s_A=4$, $s_B=6$, $s_D=-3$, $s_E=-7$):
\begin{center}
\small
\begin{tabular}{@{}cccccc@{}}
\toprule
\textbf{Nodo} & $A$ & $B$ & $C$ & $D$ & $E$\\
\midrule
$e_i = s_i + \text{entrate} - \text{uscite}$ & $4+0-0$ & $6+0-0$ & $0+0-0$ & $-3+0-3$ & $-7+3-0$\\
\textbf{sbilanciamento} & $+4$ & $+6$ & $0$ & $-6$ & $-4$\\
\bottomrule
\end{tabular}
\end{center}
Eccessi totali $=10=$ difetti totali. In $G_x$ l'arco $D\to E$ compare solo all'indietro
($E\to D$, costo $+3$): nessun costo residuo negativo, quindi lo pseudoflusso è minimale. \checkmark

\medskip
\hrule
\medskip

% ==================== ITERAZIONE 1 ====================
\textbf{Iterazione 1 --- coppia $s=A$ ($e_A=+4$), $t=E$ ($e_E=-4$).}

Cammini $A\rightsquigarrow E$ in $G_x$ e loro costo:
$A\!\to\!C\!\to\!E = 1+2 = \mathbf{3}$ \ (minimo), \quad
$A\!\to\!D$ non raggiunge $E$ (l'arco $D\to E$ è saturo).

\begin{center}
\small
\begin{tabular}{@{}ccccc@{}}
\toprule
\textbf{Arco residuo} & \textbf{Tipo} & \textbf{Formula residuo $r$} & \textbf{Residuo num.} & \textbf{Costo residuo}\\
\midrule
$A \to C$ & Diretto & $u_{AC}-x_{AC} = 4-0$ & $4$ & $+1$\\
$C \to E$ & Diretto & $u_{CE}-x_{CE} = 5-0$ & $5$ & $+2$\\
\bottomrule
\end{tabular}
\end{center}
\begin{itemize}[itemsep=0.15em]
\item \textbf{Passo}: $\delta = \min(e_A,\,|e_E|,\,4,\,5) = \min(4,4,4,5) = \mathbf{4}$.
\item \textbf{Aggiornamento}: $x_{AC}\leftarrow \mathbf{4}$ (saturo), \quad $x_{CE}\leftarrow \mathbf{4}$.
\item \textbf{Nuovi sbilanciamenti}: $e_A = 4-4 = \mathbf{0}$, \quad $e_E = -4+4 = \mathbf{0}$.
\item \textbf{Costo}: $z_1 = z_0 + \delta\cdot c(\text{cammino}) = -9 + 4\cdot 3 = \mathbf{3}$.
\end{itemize}

\begin{center}
\begin{tikzpicture}[>=Latex,scale=0.85,transform shape]
\mcfnodes
\ra{->}{A}{D}{+6}
\ra{->}{A}{C}{+1}
\ra{->}{C}{D}{+4}
\ra{->}{B}{A}{+1}
\ra{->}{B}{C}{+1}
\ra{->}{C}{E}{+2}
\ra{->}{B}{E}{+6}
\ra{<-}{D}{E}{-3}
\cy{A}{C}{+1}
\cy{C}{E}{+2}
\end{tikzpicture}
\par\vspace{1mm}{\scriptsize $G_x$ iniziale: evidenziato il cammino minimo $A\to C\to E$ ($c=3$, $\delta=4$)}
\end{center}

\medskip
\hrule
\medskip

% ==================== ITERAZIONE 2 ====================
\textbf{Iterazione 2 --- coppia $s=B$ ($e_B=+6$), $t=D$ ($e_D=-6$).}

Ora $A\to C$ è saturo (in $G_x$ resta solo $C\to A$ con costo $-1$). Cammini $B\rightsquigarrow D$:
\begin{gather*}
B\!\to\!C\!\to\!D = 1+4 = \mathbf{5}\ (\text{minimo}),\qquad
B\!\to\!C\!\to\!E\!\to\!D = 1+2+3 = 6,\\
B\!\to\!A\!\to\!D = 1+6 = 7,\qquad
B\!\to\!E\!\to\!D = 6+3 = 9
\end{gather*}
\begin{center}
\small
\begin{tabular}{@{}ccccc@{}}
\toprule
\textbf{Arco residuo} & \textbf{Tipo} & \textbf{Formula residuo $r$} & \textbf{Residuo num.} & \textbf{Costo residuo}\\
\midrule
$B \to C$ & Diretto & $u_{BC}-x_{BC} = 4-0$ & $4$ & $+1$\\
$C \to D$ & Diretto & $u_{CD}-x_{CD} = 4-0$ & $4$ & $+4$\\
\bottomrule
\end{tabular}
\end{center}
\begin{itemize}[itemsep=0.15em]
\item \textbf{Passo}: $\delta = \min(6,\,6,\,4,\,4) = \mathbf{4}$.
\item \textbf{Aggiornamento}: $x_{BC}\leftarrow \mathbf{4}$ (saturo), \quad $x_{CD}\leftarrow \mathbf{4}$ (saturo).
\item \textbf{Nuovi sbilanciamenti}: $e_B = 6-4 = \mathbf{+2}$, \quad $e_D = -6+4 = \mathbf{-2}$.
\item \textbf{Costo}: $z_2 = 3 + 4\cdot 5 = \mathbf{23}$.
\end{itemize}

\begin{center}
\begin{tikzpicture}[>=Latex,scale=0.85,transform shape]
\mcfnodes
\ra{->}{A}{D}{+6}
\ra{<-}{A}{C}{+1}
\ra{->}{C}{D}{+4}
\ra{->}{B}{A}{+1}
\ra{->}{B}{C}{+1}
\ra{<->}{C}{E}{+2}
\ra{->}{B}{E}{+6}
\ra{<-}{D}{E}{-3}
\cy{B}{C}{+1}
\cy{C}{D}{+4}
\end{tikzpicture}
\par\vspace{1mm}{\scriptsize $G_x$ post it. 1: evidenziato $B\to C\to D$ ($c=5$, $\delta=4$)}
\end{center}

\medskip
\hrule
\medskip

% ==================== ITERAZIONE 3 ====================
\textbf{Iterazione 3 --- coppia $s=B$ ($e_B=+2$), $t=D$ ($e_D=-2$).}

$B\to C$ e $C\to D$ sono ora saturi: da $B$ restano solo $B\to A$ e $B\to E$.
\[
B\!\to\!A\!\to\!D = 1+6 = \mathbf{7}\ (\text{minimo}),\qquad B\!\to\!E\!\to\!D = 6+3 = 9
\]
\begin{center}
\small
\begin{tabular}{@{}ccccc@{}}
\toprule
\textbf{Arco residuo} & \textbf{Tipo} & \textbf{Formula residuo $r$} & \textbf{Residuo num.} & \textbf{Costo residuo}\\
\midrule
$B \to A$ & Diretto & $u_{BA}-x_{BA} = 5-0$ & $5$ & $+1$\\
$A \to D$ & Diretto & $u_{AD}-x_{AD} = 6-0$ & $6$ & $+6$\\
\bottomrule
\end{tabular}
\end{center}
\begin{itemize}[itemsep=0.15em]
\item \textbf{Passo}: $\delta = \min(2,\,2,\,5,\,6) = \mathbf{2}$.
\item \textbf{Aggiornamento}: $x_{BA}\leftarrow \mathbf{2}$, \quad $x_{AD}\leftarrow \mathbf{2}$.
\item \textbf{Nuovi sbilanciamenti}: $e_B = \mathbf{0}$, \quad $e_D = \mathbf{0}$: tutti i nodi sono bilanciati, \textbf{stop}.
\item \textbf{Costo}: $z_3 = 23 + 2\cdot 7 = \mathbf{37}$.
\end{itemize}

\begin{center}
\begin{minipage}[t]{0.48\textwidth}\centering
\begin{tikzpicture}[>=Latex,scale=0.75,transform shape]
\mcfnodes
\ra{->}{A}{D}{+6}
\ra{<-}{A}{C}{+1}
\ra{<-}{C}{D}{+4}
\ra{->}{B}{A}{+1}
\ra{<-}{B}{C}{+1}
\ra{<->}{C}{E}{+2}
\ra{->}{B}{E}{+6}
\ra{<-}{D}{E}{-3}
\cy{B}{A}{+1}
\cy{A}{D}{+6}
\end{tikzpicture}
\par\vspace{1mm}{\scriptsize $G_x$ post it. 2: evidenziato $B\to A\to D$ ($c=7$, $\delta=2$)}
\end{minipage}\hfill
\begin{minipage}[t]{0.48\textwidth}\centering
\begin{tikzpicture}[>=Latex,scale=0.75,transform shape]
\mcfnodes
\fa{A}{D}{6,\,6}{2}
\fa{A}{C}{4,\,1}{4}
\fa{C}{D}{4,\,4}{4}
\fa{B}{A}{5,\,1}{2}
\fa{B}{C}{4,\,1}{4}
\fa{C}{E}{5,\,2}{4}
\fa{B}{E}{3,\,6}{0}
\fa{D}{E}{3,\,-3}{3}
\end{tikzpicture}
\par\vspace{1mm}{\scriptsize Flusso ottimo $x^*$, costo minimo $\mathbf{37}$}
\end{minipage}
\end{center}

\medskip
\hrule
\medskip

\textbf{Soluzione ottima finale} (identica a quella trovata per cancellazione dei cicli):
\[
x_{AC}^*=4,\quad x_{AD}^*=2,\quad x_{BA}^*=2,\quad x_{BC}^*=4,\quad x_{BE}^*=0,
\quad x_{CD}^*=4,\quad x_{CE}^*=4,\quad x_{DE}^*=3
\]
\[
z^* = 4\cdot 1 + 2\cdot 6 + 2\cdot 1 + 4\cdot 1 + 0\cdot 6 + 4\cdot 4 + 4\cdot 2 + 3\cdot(-3) = \mathbf{37}
\]
\textbf{Confronto fra i due algoritmi.} La cancellazione dei cicli mantiene sempre
l'\emph{ammissibilità} (parte da un flusso di costo 48 e scende a 37); i cammini minimi
successivi mantengono sempre l'\emph{ottimalità nei costi} (parte da $-9$ e sale a 37,
riducendo lo sbilanciamento a ogni passo). Il valore finale coincide.
\end{es}

\begin{teo}{pseudoflussi, cammini aumentanti e ottimalità}
\textbf{Pseudoflusso.} Un vettore $x$ con $0\le x_{ij}\le u_{ij}$ per ogni arco,
senza vincolo di conservazione. Ogni nodo ha uno \emph{sbilanciamento} (eccesso o difetto);
lo pseudoflusso è \textbf{ammissibile} (cioè è un flusso) se e solo se tutti gli
sbilanciamenti sono nulli.

\textbf{Pseudoflusso minimale.} Uno pseudoflusso il cui grafo residuo $G_x$ non contiene
cicli di costo negativo: è di costo minimo fra tutti gli pseudoflussi con gli stessi
sbilanciamenti.

\textbf{Condizione di ottimalità.} Un flusso ammissibile $x$ è di costo minimo se e solo se
$G_x$ non contiene cicli di costo negativo.

\textbf{Trasformazione fra pseudoflussi.} Due pseudoflussi qualunque con gli stessi
sbilanciamenti differiscono per una combinazione di cicli di $G_x$: si passa dall'uno all'altro
inviando flusso lungo cicli, e la differenza di costo è la somma dei costi di quei cicli.
Se invece gli sbilanciamenti sono diversi, la differenza è una combinazione di cicli
\emph{e} di cammini fra nodi sbilanciati.

\textbf{Effetto di un cammino aumentante di costo minimo.} Spingere flusso lungo un cammino
di costo minimo da un nodo in eccesso a uno in difetto trasforma uno pseudoflusso minimale
in un altro pseudoflusso ancora minimale, riducendo lo sbilanciamento complessivo.
Da qui la correttezza dei cammini minimi successivi.

\textbf{Simmetria dei due algoritmi.} La cancellazione dei cicli mantiene
l'\emph{ammissibilità} e migliora il costo (algoritmo primale); i cammini minimi successivi
mantengono l'\emph{ottimalità dei costi} e ripristinano l'ammissibilità (algoritmo duale).
\end{teo}

\newpage
\subsection{Matching di massima cardinalità}

Grafo bipartito con lati $U$ e $V$ e insieme di archi $E$.

\begin{ric}{matching risolto come flusso massimo}
\begin{enumerate}
\item Aggiungi una sorgente $s$ e un arco $s \to u$ di capacità 1 per ogni $u \in U$.
\item Orienta ogni arco di $E$ come $u \to v$, con capacità 1 (o $+\infty$, è indifferente:
      i colli di bottiglia sono agli estremi).
\item Aggiungi un pozzo $t$ e un arco $v \to t$ di capacità 1 per ogni $v \in V$.
\item Esegui Edmonds--Karp sulla rete così costruita.
\item Il valore del flusso massimo è la cardinalità massima del matching;
      gli archi $u \to v$ con flusso 1 sono le coppie accoppiate.
\end{enumerate}
\end{ric}

\begin{es}{matching con una riassegnazione}
$U = \{u_1,u_2,u_3\}$, $V = \{v_1,v_2,v_3\}$, archi
$E = \{u_1v_1,\ u_1v_2,\ u_2v_1,\ u_3v_2,\ u_3v_3\}$.

\begin{center}
\begin{tikzpicture}[>=Latex,every node/.style={font=\small},
  nd/.style={circle,draw=navy,thick,minimum size=7mm,inner sep=1pt}]
\node[nd] (u1) at (0,2) {$u_1$};
\node[nd] (u2) at (0,1) {$u_2$};
\node[nd] (u3) at (0,0) {$u_3$};
\node[nd] (v1) at (3,2) {$v_1$};
\node[nd] (v2) at (3,1) {$v_2$};
\node[nd] (v3) at (3,0) {$v_3$};
\node[nd] (s) at (-2,1) {$s$};
\node[nd] (t) at (5,1) {$t$};
\draw[->,thick] (s) -- (u1); \draw[->,thick] (s) -- (u2); \draw[->,thick] (s) -- (u3);
\draw[->,thick] (v1) -- (t); \draw[->,thick] (v2) -- (t); \draw[->,thick] (v3) -- (t);
\draw[->,thick,teal] (u1) -- (v1);
\draw[->,thick,teal] (u1) -- (v2);
\draw[->,thick,teal] (u2) -- (v1);
\draw[->,thick,teal] (u3) -- (v2);
\draw[->,thick,teal] (u3) -- (v3);
\end{tikzpicture}
\end{center}
Tutte le capacità valgono 1. Ordine BFS: per indice crescente.

\textbf{Iterazione 1.} Cammino più corto: $s \to u_1 \to v_1 \to t$ (3 archi). $\delta = 1$,
flusso $=1$. Matching parziale: $u_1$--$v_1$.

\textbf{Iterazione 2.} $s\to u_2 \to v_1$ è bloccato perché $v_1\to t$ è saturo;
il cammino più corto disponibile è $s \to u_3 \to v_2 \to t$ (3 archi). $\delta=1$, flusso $=2$.
Matching: $u_1$--$v_1$, $u_3$--$v_2$.

\textbf{Iterazione 3.} Resta libero solo $u_2$, che vede solo $v_1$, già occupato.
BFS in $G_x$: $s \to u_2 \to v_1 \to u_1$ (\emph{arco inverso}: si disfa $u_1$--$v_1$)
$\to v_2 \to u_3$ (\emph{arco inverso}: si disfa $u_3$--$v_2$) $\to v_3 \to t$.
Sette archi, $\delta = 1$, flusso $= 3$.

\textbf{Iterazione 4.} Da $s$ nessun arco residuo uscente ($s\to u_1,u_2,u_3$ tutti saturi):
\textbf{stop}.

\medskip
\textbf{Soluzione.} Matching massimo di cardinalità \textbf{3}:
\[ u_1\text{--}v_2, \qquad u_2\text{--}v_1, \qquad u_3\text{--}v_3 \]
Si noti come i due archi inversi dell'iterazione 3 abbiano \emph{riassegnato} $u_1$ e $u_3$
per fare posto a $u_2$: è esattamente il meccanismo del cammino aumentante.

\textbf{Certificato.} $S=\{s\}$ non funziona qui perché il flusso vale 3 = numero di archi
uscenti da $s$: il taglio $(\{s\},\text{resto})$ ha capacità 3 = valore del flusso, quindi è
minimo e certifica l'ottimalità (nessun matching può superare $|U| = 3$).
\end{es}

L'ottimalità è sempre certificata dal taglio minimo. È il \textbf{teorema di König}:
in un grafo bipartito la cardinalità massima di un matching è uguale alla
cardinalità minima di una copertura per vertici; la copertura minima si legge dal taglio
minimo come $(U \cap T) \cup (V \cap S)$.

\newpage
%================================================================
\section{Le domande di teoria che possono uscire}

Elenco ufficiale dei possibili approfondimenti teorici (a.a.\ 2025--26).
Le voci in \textbf{grassetto} sono già state chieste nei compiti recenti.

\textbf{Ottimizzazione lineare}
\begin{itemize}[itemsep=0.15em]
\item Descrivere i quattro casi di soluzione di un problema di ottimizzazione
      (ottimo finito, ottimi multipli, illimitato, vuoto).
\item Definire problema di ottimizzazione, di esistenza e di certificato.
\item Definire un algoritmo di tipo primale e duale e fornirne un esempio
      (cancellazione cicli / cammini minimi successivi).
\item Definire un rilassamento di un problema di ottimizzazione e illustrarne un esempio
      (rilassamento continuo nel Branch-and-Bound).
\item Modellazione con vincoli lineari di relazioni logiche (almeno un caso):
      implicazione $x_b \le x_a$, incompatibilità $x_h + x_k \le 1$, big-M.
\item Modellazione con vincoli lineari di funzioni lineari a tratti e convesse.
\item \textbf{Definire una soluzione base ed una soluzione base ammissibile di un problema PL.}
\item Enunciare il teorema di Motzkin (decomposizione) per un poliedro.
\item Enunciare e dimostrare la proprietà fondamentale della PL.
\item Definire il problema primale e duale e fornire un esempio.
\item \textbf{Discutere l'interpretazione delle variabili duali} (prezzi ombra).
\item \textbf{Enunciare e dimostrare il teorema debole della dualità.}
\item \textbf{Definire una direzione ammissibile e di crescita.}
\item Enunciare la condizione di ottimalità per la PL.
\end{itemize}

\textbf{Ottimizzazione su rete}
\begin{itemize}[itemsep=0.15em]
\item Definire il problema del flusso massimo su una rete ed il relativo modello matematico.
\item Definire il problema del flusso di costo minimo su una rete ed il relativo modello matematico.
\item Data una rete definire un taglio $s$-$t$ ed illustrarne le proprietà.
\item Data una rete ed un flusso definire il grafo residuo corrispondente.
\item Definire un cammino aumentante per il flusso massimo sul grafo residuo.
\item Discutere la correttezza dell'algoritmo di Ford--Fulkerson.
\item \textbf{Enunciare e dimostrare il teorema max-flow-min-cut.}
\item \textbf{Definire uno pseudoflusso e discutere le condizioni della sua ammissibilità.}
\item Definire un cammino aumentante per il flusso a costo minimo sul grafo residuo.
\item Definire uno pseudoflusso minimale.
\item Discutere l'effetto di un cammino aumentante a costo minimo su uno pseudoflusso minimale.
\item Discutere come ottenere un flusso ammissibile, se esiste, per un problema di flusso
      a costo minimo (rete ampliata con $s$ e $t$ fittizi + flusso massimo).
\item \textbf{In cosa differisce Edmonds--Karp dalle altre implementazioni per il flusso massimo}
      (visita in ampiezza, cammino con meno archi, complessità indipendente dalle capacità).
\item \textbf{Definire un flusso ammissibile.}
\end{itemize}

\textbf{I quattro casi di soluzione} (domanda ricorrente, risposta pronta):
\begin{enumerate}[itemsep=0.15em]
\item \textbf{Ottimo finito unico}: esiste una sola soluzione ottima.
\item \textbf{Ottimi multipli}: l'insieme delle soluzioni ottime è una faccia del poliedro
      (spigolo o faccia intera); segnalato da una variabile duale nulla in base.
\item \textbf{Illimitato}: esiste una direzione ammissibile di crescita infinita;
      il duale è vuoto.
\item \textbf{Vuoto (inammissibile)}: nessuna soluzione soddisfa i vincoli.
\end{enumerate}

\newpage
%================================================================
\section{Checklist prima di consegnare}

\begin{idea}
\textbf{Modellazione}
\begin{itemize}
\item Tutte le variabili definite a parole, con il loro dominio nella riga finale.
\item Ogni vincolo quantificato con $\forall i$ oppure $\forall j$.
\item Costi fissi legati alle variabili di attivazione ($\le b_i y_i$ oppure $x_{ij}\le y_i$).
\item Riletto il testo frase per frase: ogni frase ha prodotto un vincolo?
\end{itemize}
\end{idea}

\begin{idea}
\textbf{Simplesso}
\begin{itemize}
\item Vincoli normalizzati in $\max c\,x$, $Ax\le b$ e \emph{numerati per colonne}.
\item $A$, $b$, $c$ scritti esplicitamente prima di iniziare.
\item Regione ammissibile colorata, vertici marcati, gradiente $c$ disegnato col verso.
\item Una tabella per ogni iterazione con $B$, $A_B^{-1}$, $x$, $y_B$, $h$, $\xi$, $\lambda_i$, $k$.
\item Ogni $A_B^{-1}$ verificata con $A_B A_B^{-1} = I$.
\item Vettore $y$ finale riportato come certificato di ottimalità, e verifica $c\,x = y\,b$.
\item Il vertice ottimo del simplesso coincide con quello letto sul grafico.
\end{itemize}
\end{idea}

\begin{idea}
\textbf{Reti}
\begin{itemize}
\item Grafo residuo ridisegnato a ogni iterazione, con il $\delta$ evidenziato.
\item Cammino o ciclo scelto scritto esplicitamente in lettere, non solo indicato sul disegno.
\item Regola di rottura dei pari nella BFS dichiarata a inizio esercizio.
\item Taglio minimo dato come coppia $(S, T)$, con i tre controlli numerici.
\item Nel costo minimo: bilanci di ogni nodo verificati alla fine, costo totale ricalcolato.
\end{itemize}
\end{idea}

\begin{idea}
\textbf{Teoria}
\begin{itemize}
\item Prima l'enunciato, poi la dimostrazione, in poche righe pulite.
\item Le definizioni vanno scritte per intero: base, base ammissibile, direzione ammissibile,
      direzione di crescita, flusso ammissibile, pseudoflusso, pseudoflusso minimale, taglio.
\item Nelle dimostrazioni, dichiarare sempre \emph{dove} si usa ogni ipotesi
      (es. «qui uso $y\ge0$ e $Ax\le b$»).
\end{itemize}
\end{idea}

\end{document}
